% SPDX-FileCopyrightText: 2026 Will Cook
% SPDX-License-Identifier: CC-BY-4.0
%
% Long-form reasoning record seeded from the complete problem note.
% The short note and this record are now independent publication units.
% This flat manuscript is generated from reviewable parts below.
\documentclass[11pt]{article}
% SPDX-FileCopyrightText: 2026 Will Cook
% SPDX-License-Identifier: CC-BY-4.0
%
% Shared preamble for the Erdős Problem Notes: the short, standalone,
% problem-owned manuscripts covering the expansion library `ErdosProblems`.
%
% One note owns one Erdős problem.  Each note repeats whatever context its own
% reader needs, so that a reader who arrives at a single problem never has to
% assemble it from the gateway paper.  Deliberate repetition across notes is the
% design, not an oversight.
%
% Source links resolve against one pinned formal-source commit, declared once
% below.  `scripts/check_problem_note_sources.py` verifies that every authored
% (file, line, declaration) triple names that exact declaration in the pinned
% snapshot, so a link cannot silently rot as later work moves lines.

\usepackage{paper-house-style}
\usepackage{array}
\usepackage{booktabs}
\usepackage{enumitem}
\setlist{topsep=0.35em,itemsep=0.2em}
\usepackage[colorlinks=true,linkcolor=linkinternal,urlcolor=linkexternal,
  citecolor=linkinternal]{hyperref}
\hypersetup{bookmarksnumbered=true,bookmarksopen=true,bookmarksopenlevel=1,
  pdfauthor={Will Cook},pdflang={en-GB}}

% ---------------------------------------------------------------------------
% Pinned formal source.  \commit names the checkpoint every link resolves
% against; the checker reads that snapshot rather than the working tree, so the
% printed coordinates stay correct however later waves move declarations.
% ---------------------------------------------------------------------------
\newcommand{\commit}{e5fd26a6d1b1a346430205eab5385b4c9565d004}
\newcommand{\commitshort}{e5fd26a6d1b1}
\newcommand{\repobase}{https://github.com/wcook04/plectis-lean-erdos249-257/blob/\commit}
\newcommand{\sourceurl}{https://github.com/wcook04/plectis-lean-erdos249-257/tree/\commit}
\newcommand{\PX}{ErdosProblems}

% Declaration names stay inside link targets.  The printed label is either a
% spaced, lower-case reading of the declaration name (\lref) or a short authored
% phrase (\lword); neither prints a module path or a raw Lean identifier.
\ExplSyntaxOn
\cs_new_protected:Npn \noteleanlabel #1
  {
    \tl_set:Nx \l_tmpa_tl { \tl_to_str:n {#1} }
    \regex_replace_all:nnN { _+ } { \c{space} } \l_tmpa_tl
    \regex_replace_all:nnN
      { ([a-z0-9])([A-Z]) }
      { \1\c{space}\2 }
      \l_tmpa_tl
    \tl_set:Nx \l_tmpa_tl { \tl_lower_case:n { \tl_use:N \l_tmpa_tl } }
    \tl_use:N \l_tmpa_tl
  }
\ExplSyntaxOff
\newcommand{\leanlabel}[1]{\noteleanlabel{#1}}
\newcommand{\lref}[3]{\href{\repobase/\PX/#1\#L#2}{\leanlabel{#3}}}
\newcommand{\lword}[4]{\href{\repobase/\PX/#1\#L#2}{#4}}
\newcommand{\lloc}[2]{\href{\repobase/\PX/#1\#L#2}{Lean source}}

% Links into a sibling library, named repository-relative.  The expansion
% library is implicit in \lref/\lword, because that is where problem-owned
% work lands.  Several problems have their headline theorem in the older
% reviewed #249/#257 corpus, and a note that could not name that library
% would have to paraphrase a checked statement instead of linking it.
\newcommand{\mref}[3]{\href{\repobase/#1\#L#2}{\leanlabel{#3}}}
\newcommand{\mword}[4]{\href{\repobase/#1\#L#2}{#4}}
\newcommand{\mloc}[2]{\href{\repobase/#1\#L#2}{Lean source}}
\emergencystretch=2em

\newtheorem{theorem}{Theorem}[section]
\newtheorem{proposition}[theorem]{Proposition}
\newtheorem{lemma}[theorem]{Lemma}
\newtheorem{corollary}[theorem]{Corollary}
\theoremstyle{definition}
\newtheorem{definition}[theorem]{Definition}
\newtheorem{example}[theorem]{Example}
\newtheorem{problem}[theorem]{Problem}
\theoremstyle{remark}
\newtheorem*{remark}{Remark}

\newcommand{\N}{\mathbb{N}}
\newcommand{\Npos}{\mathbb{N}_{>0}}
\newcommand{\Z}{\mathbb{Z}}
\newcommand{\Q}{\mathbb{Q}}
\newcommand{\R}{\mathbb{R}}
\newcommand{\lcm}{\operatorname{lcm}}
\newcommand{\ph}{\varphi}

% The series line printed under every note's title block.
\newcommand{\seriesline}{Erd\H{o}s Problem Notes}

% Production provenance.  The wording is fixed by the corpus provenance
% contract and reproduced verbatim in every manuscript in this repository, so
% the papers cannot drift into different accounts of how they were made.
\newcommand{\noteprovenance}{\provenancenote{\emph{Authorship and AI use.} The
prose, formal proofs, and software in this version were drafted and revised by
large-language-model agents under Will Cook's direction.  Cook set the
objectives and acceptance criteria, reviewed and authorised the manuscript, and
is responsible for its contents.  The agents are production tools, not authors.
See \emph{Statements and declarations}.}}

% The status paragraph every note carries, in its own words, before any result.
% Kernel checking establishes the exact Lean proposition; it does not establish
% that the proposition is the interesting one, that it is new, or that the
% Erdős problem is closed.
\newcommand{\statusboundary}{%
  \emph{Status.}  The problem treated here is open, and this note does not
  close it.  Every statement below marked as checked is a proposition that the
  pinned Lean kernel accepts from the sources this note links to, with no
  \texttt{sorry}, no added axiom, and no unchecked evaluation.  That is a claim
  about the formal statement, not about its mathematical interest, its
  novelty, or the original problem.  The unresolved obligations are named
  exactly, in their own section, and none of the finite computations,
  reductions, or no-go results here removes one of them.}

\renewcommand{\commit}{92b88dc1bbe099aa73bcc900c3c405e9ba5c2334}
\renewcommand{\commitshort}{92b88dc1bbe0}
\usepackage{longtable}
\hypersetup{
  pdftitle={Denominators and rationality criteria for the factorial-denominator series (Erdős Problem 68)},
  pdfsubject={Erdős Problem Notes: Problem 68},
  pdfkeywords={irrationality, factorial denominators, factorial digits, prime-power cancellation, support bounds, Lean 4},
}

% Source link targets retain the original manuscript commits and paths.
\newcommand{\Ser}{S}
\newcommand{\chn}[2]{V_{#1}(#2)}
\newcommand{\Lc}{L}
\newcommand{\den}{\operatorname{den}}

\runninghead{Erd\H{o}s \#68: denominators and rationality criteria}
\title{Denominators and Rationality Criteria\\for $\sum_{n\ge2}(n!-1)^{-1}$}
\subtitle{Detailed proofs and further deductions for Erd\H{o}s Problem~\#68}
\author{Will Cook}
\date{1 September 2026; revised 18 September 2026}

\setlength{\emergencystretch}{3em}
% Permit breaks inside long declaration names in the generated catalogue.
\usepackage{xurl}
\begin{document}
% ---- part core ----
% , - part core , -
\maketitle
\noteprovenance

\begin{abstract}
Let $S=\sum_{n\ge2}(n!-1)^{-1}$ and
$\Lc_N=\lcm(2!-1,\ldots,N!-1)$. We prove
\[
 \liminf_{N\to\infty}\frac{\log\Lc_N}{N^{3/2}\log N}
 \ge\frac{2\sqrt2}{3}.
\]
The proof compares the product of a terminal block of denominators with
their pairwise gcds, using
$\gcd(i!-1,j!-1)\mid j!/i!-1$ rather than an external multiplicity theorem. This is a bound before reduction. To study
the reduced partial sums, we specialise a valuation formula of Louwsma
and Martino to reciprocal sums. Exact examples show that the primes $139$ and
$2593$ are absent from the reduced denominators of every partial sum from
indices $137$ and $2591$ onward, respectively.

The coefficient constructions give integer linear forms $MS+k$.
We classify the vectors that cancel prescribed initial weighted sums,
determine their attainable moments $M$ on indices at least $2$, and give a
lower bound for their largest support index when the moment obeys a stated
upper bound there. At fixed moment the remainder changes only by an integer.
The rationality criteria use factorial digits and the next integer above
a scaled partial sum. The digit argument also applies to
$\sum_{n\ge2}1/(n!+t)$ for integers $t\ge-1$, recovering the irrationality
of $e$ at $t=0$.

Finally, two exact finite calculations give
$q\nmid299999!$ and $q\ge2^{39990}>10^{12038}$ for any rational
representation $S=a/q$, $q>0$. The full carry calculation and the
continued-fraction calculation are external to Lean. These exclusions do not prove irrationality; the additional tail and
residue inequalities that would suffice are stated separately.
\end{abstract}

\paragraph{Organisation.}
The first two sections keep two different questions separate: which prime
powers survive reduction (\S\ref{long68:sec:prime-powers}), and how large a
common denominator must be before reduction (\S\ref{long68:sec:lcm}).
The tail comparisons in \S\ref{long68:sec:residue} require information not
supplied by either answer. For the finite exclusions, only the carry
criterion in \S\ref{long68:sec:carry} and the computations in
\S\ref{long68:sec:finite} are needed. The coefficient constructions in
\S\ref{long68:sec:channels} are a separate approach; \S\ref{long68:sec:open}
identifies the remaining inequalities, and the appendices give further
deductions, unsuccessful approaches, and formal-source links.

\section{Which prime powers survive reduction}\label{long68:sec:prime-powers}

\begin{problem}[Erd\H{o}s \#68]\label{long68:res:problem}
Is $\Ser=\sum_{n\ge2}(n!-1)^{-1}$ irrational?
\end{problem}

\noindent Erd\H{o}s states the question on p.~102 of his 1988 survey and, in
the same passage, records the expectation that $\sum_n1/(n!+t)$ is
irrational, indeed transcendental, for every integer $t$
\cite[p.~102]{erdos1988}. Here the sum starts sufficiently far out that
$n!+t>0$; changing that starting index adds or removes only a rational
number. In particular, no term with zero denominator is included. The
expectation for all integer shifts is not proved here. Numbering follows
\href{https://www.erdosproblems.com/68}{Bloom's Erd\H{o}s problem catalogue}
\cite{bloom}.

Put
\[
 d_n=n!-1,\qquad H_M=\sum_{n=2}^{M}\frac1{d_n},\qquad
 \Lc_M=\lcm(d_2,\ldots,d_M),\qquad
 A_M=\sum_{n=2}^{M}\frac{\Lc_M}{d_n},
\]
so that $H_M=A_M/\Lc_M$.  Throughout, $\den(x)$ is the positive denominator
of a rational number in lowest terms, and $v_p(a)$ is the exponent of the
prime $p$ in a positive integer $a$. The first question is which prime
power present in $\Lc_M$ remains in $\den(H_M)$.  For the largest exponent the
answer follows by specialising Louwsma and Martino's valuation formula
for an elementary symmetric sum \cite[Lemma~4.1, p.~10]{louwsma-martino}.
Indeed, $\sum_i1/x_i=(\sum_i\prod_{j\ne i}x_j)/\prod_i x_i$ for positive
integers $x_i$. Their lemma evaluates the numerator's valuation; subtracting
the product's valuation gives the reciprocal-sum form. The proof below
uses the least common multiple instead of the product.

The need to check cancellation is already visible in $1/3+1/15=2/5$:
the common denominator contains $3$, but the reduced denominator does not.
By contrast, $1/9+1/45=2/15$ only lowers the exponent of $3$ from $2$ to $1$.
The theorem below tests whether the largest exponent is preserved; a
failure of that test need not remove the prime completely. The formula
is standard, whereas the two factorial-prefix examples below and the later
tail comparison concern this particular series.

\label{long68:res:lead-prime-pole}\phantomsection
\begin{theorem}[maximal prime-power survival]\label{long68:res:prime-pole}
Let $M\ge2$, let $p$ be a prime dividing $\Lc_M$, and put $e=v_p(\Lc_M)$.
Let $J=\{n:2\le n\le M,\ v_p(d_n)=e\}$ and write $d_n=p^eu_n$ for $n\in J$.
Then, with inverses in $\mathbb F_p$,
\begin{equation}\label{long68:eq:prime-pole-survival}
 v_p\bigl(\den(H_M)\bigr)=e
 \quad\Longleftrightarrow\quad
 \sum_{n\in J}u_n^{-1}\ne0\quad\hbox{in }\mathbb F_p.
\end{equation}
\end{theorem}
\leannote{Lean: \lproof{ErdosProblems/Erdos68/PaperCompletePrimePole.lean}{117}{maximal_prime_power_survival}.}

\begin{proof}
Write $\Lc_M=p^eW$ with $p\nmid W$.  If $v_p(d_n)<e$, then $p\mid\Lc_M/d_n$.
If $n\in J$, then $u_n(\Lc_M/d_n)=W$, so $\Lc_M/d_n\equiv Wu_n^{-1}\pmod p$.
Summing over $2\le n\le M$ gives
\begin{equation}\label{long68:eq:prime-pole-residue}
 A_M\equiv \frac{\Lc_M}{p^e}\sum_{n\in J}u_n^{-1}\pmod p.
\end{equation}


Finally $\den(H_M)=\Lc_M/\gcd(A_M,\Lc_M)$, whose $p$-valuation is $e$ exactly
when $p\nmid A_M$.  Since $W$ is invertible modulo $p$, that is exactly the
stated condition.
\end{proof}

The complete valuation identity is
\[
 v_p\bigl(\den(H_M)\bigr)=\max\{0,e-v_p(A_M)\}.
\]
A zero residue in \eqref{long68:eq:prime-pole-survival} therefore lowers
the exponent, whereas complete cancellation requires $v_p(A_M)\ge e$.
For $e>1$, reduction modulo $p$ alone cannot distinguish partial from
complete cancellation.  The proof uses no property of factorials beyond the
displayed denominators, so it applies to any finite sum of reciprocals of positive integers.
For $d_n=n!-1$, the examples below show that complete cancellation can
actually occur; the general valuation formula alone does not predict it.

\paragraph{Two complete cancellations.}

A unique maximal exponent makes the residue sum a single nonzero term,
so the condition holds automatically. With several maximal terms it can
fail. The following examples show cancellation in the factorial sequence
itself, not just in arbitrary rational sums.
Take $M=p-1$. The table lists every index with $p\mid d_n$, and every
one has valuation exactly one.
\begin{center}
\small
\begin{tabular}{rlll}
\toprule
\papertablehead
$p$ & indices $n$ & $d_n/p\pmod p$ & inverses modulo $p$\\
\midrule
$139$ & $69,\ 122,\ 137$ & $6,\ 49,\ 73$ & $116,\ 122,\ 40$\\
$2593$ & $349,\ 2243,\ 2591$ & $1508,\ 1566,\ 1678$ & $1367,\ 356,\ 870$\\
\bottomrule
\end{tabular}
\end{center}
The inverse sums are $278=2\cdot139$ and $2593$.  By
Theorem~\ref{long68:res:prime-pole}, $139\nmid\den(H_{138})$ and
$2593\nmid\den(H_{2592})$: each prime divides the common denominator of its
prefix and vanishes from the reduced one.  Both hit lists and both valuations
are checked by the recurrence
\[
 r_1=1,\qquad r_n\equiv nr_{n-1}\pmod{p^2},\qquad 0\le r_n<p^2,
\]
run through $n=p-1$.  A hit is an index with $r_n\equiv1\pmod p$, its lifted
cofactor is $(r_n-1)/p$ modulo $p$. Valuation at least two would give
$r_n=1$, which occurs for neither prime among the tested indices
$2\le n\le p-1$.  That enumeration is a finite calculation
separate from the theorem.

The recurrence is run modulo $p^2$, not merely modulo $p$: reduction
modulo $p$ would locate hits but could not establish that their valuations
are exactly one. This is why the zero inverse sums prove complete
cancellation in these examples.

For each prime $p\ge3$, the exponent $v_p(\den(H_M))$ is constant for
$M\ge\max(2,p-2)$. Wilson's theorem gives
$(p-1)!-1\equiv-2\pmod p$, and $n!-1\equiv-1\pmod p$ for $n\ge p$.
Hence every summand after index $p-2$ has denominator coprime to $p$. If a reduced fraction
$a/b$ is followed by a summand $1/d$ with $p\nmid d$, their sum has
numerator $ad+b$ over $bd$. When $p\mid b$, this numerator is nonzero
modulo $p$, so the exponent of $p$ in the reduced denominator stays the
same. When $p\nmid b$, the new denominator remains coprime to $p$.
Consequently, the two cancellations persist in every later partial sum:
\[
 139\nmid\den(H_M)\quad(M\ge137),\qquad
 2593\nmid\den(H_M)\quad(M\ge2591).
\]
The primes still divide the corresponding common denominators $L_M$.
These are assertions about the partial sums; they do not constrain a
hypothetical denominator of their real limit $S$ without a further
approximation estimate.

We next ask where a prime first divides a denominator $d_m$. This means
it divides none of $d_2,\ldots,d_{m-1}$; it may still divide a later
denominator, so uniqueness at its first occurrence gives no uniqueness
in every larger block.

\begin{proposition}[cofinal first prime occurrences]\label{long68:res:wilson-cofinality}
For every integer $B\ge0$ there are a prime $q$ and an integer $m>B$ with
$m<q$, $q\mid m!-1$ and $\gcd(q,k!-1)=1$ for every $k$ with $2\le k<m$.
\end{proposition}
\leannote{Lean: \lproof{ErdosProblems/Erdos68/PaperCompleteExisting.lean}{180}{cofinal_first_prime_occurrences}.}

\begin{proof}
Choose a prime $q\ge B!+5$.  Wilson's theorem gives $(q-2)!\equiv1\pmod q$,
so there is a least index $m\ge2$ with $q\mid m!-1$, and $m\le q-2$.
If $m\le B$, then
$q\le m!-1\le B!-1$, contradicting the choice of $q$.  Hence $m>B$, and
minimality gives the coprimality assertions.
\end{proof}

The construction proves cofinality of first occurrences, but not the
inequality involving $q$ needed for the additional two-factor test in
\S\ref{long68:sec:residue}. The sufficient tail inequality there does not
require choosing a prime $q$. Wilson reflection limits what can be inferred
from a large prime factor alone: for odd $2\le n<q$ with $q\mid n!-1$,
the identity $(q-1-n)!\,n!\equiv(-1)^{n+1}\pmod q$ gives the reflected
congruence at $q-n-1$. It gives a second summand denominator only when
$q\ne2n+1$ and $n\le q-3$. At the Wilson endpoint $n=q-2$, it instead
returns index $1$, where $1!-1=0$ is not a denominator of this series.  The reflection identity is classical; see Stewart
\cite[p.~462, (4)]{stewart2004}.

\statusboundary

\section{The growth of the common denominator}\label{long68:sec:lcm}

Theorem~\ref{long68:res:prime-pole} concerns the denominator after reduction.  The
common denominator before reduction admits an unconditional lower bound of
its own, by an elementary argument.

\begin{lemma}[product, least common multiple, pairwise gcd]
\label{long68:res:product-lcm}
For positive integers $x_1,\ldots,x_k$,
\[
 \prod_{i=1}^{k}x_i\ \Big|\ \lcm(x_1,\ldots,x_k)\prod_{i<j}\gcd(x_i,x_j).
\]
\end{lemma}
\leannote{Lean: \lproof{ErdosProblems/Erdos68/PaperCompleteExisting.lean}{189}{product_lcm_pairwise_gcd}.}

\begin{proof}
Fix a prime $r$ and relabel so that $a_1\le\cdots\le a_k$, where
$a_i=v_r(x_i)$.  The right-hand side has $r$-valuation
$a_k+\sum_{i<j}\min(a_i,a_j)=a_k+\sum_{i=1}^{k-1}(k-i)a_i$, and the left-hand
side has $\sum_{i=1}^{k}a_i$.  The difference is
$\sum_{i=1}^{k-1}(k-i-1)a_i$, which is nonnegative.
\end{proof}

\noindent The next divisibility is the case $P=-1$ of the relation
$\gcd(i!+P(i),j!+P(j))\mid (j!/i!)P(i)-P(j)$ for $P\in\Z[X]$, obtained by
multiplying $i!+P(i)$ by $j!/i!$ and subtracting $j!+P(j)$.  Luca and
Shparlinski use this relation to bound a common divisor of $n!+P(n)$ and
$(n+h)!+P(n+h)$ in the proof of their Lemma~5~\cite{luca-shparlinski}, and Lai
uses it for a common prime-power divisor
\cite[proof of Lemma~2.4, display~(2.5)]{lai}.

The earlier spacing method of Erd\H{o}s and Stewart
\cite[\S3, pp.~516--517]{erdos-stewart1976} is related background.
We now state only the subtraction needed for the lcm argument, rather than
importing the prime-factor estimates of those papers.

\begin{lemma}[factorial-gap gcd]\label{long68:res:gap-gcd}
For $2\le i<j$, the integer $g=\gcd(i!-1,j!-1)$ divides $j!/i!-1$, and
$g\le j!/i!-1<j^{\,j-i}$.
\end{lemma}
\leannote{Lean: \lproof{ErdosProblems/Erdos68/ChannelIntegralCongruence.lean}{288}{factorial_gap_gcd_exact}.}

\begin{proof}
Both $i!$ and $j!$ are congruent to $1$ modulo $g$.  The quotient
$j!/i!=(i+1)(i+2)\cdots j$ is an integer, and $j!=i!\cdot(j!/i!)$, so
$j!/i!\equiv1\pmod g$.  Since $j!/i!\ge i+1>1$, the positive integer
$j!/i!-1$ is a multiple of $g$, whence $g\le j!/i!-1$.  Finally
$j!/i!$ is a product of $j-i$ integers each at most $j$.
\end{proof}

\begin{lemma}[segment inequality]\label{long68:res:segment}
For $2\le k\le N-1$,
\begin{equation}\label{long68:eq:segment}
 \sum_{n=N-k+1}^{N}\log(n!-1)
 \ \le\ \log \Lc_N+\binom{k+1}{3}\log N.
\end{equation}


\end{lemma}
\leannote{Lean: \lproof{ErdosProblems/Erdos68/ChannelIntegralCongruence.lean}{566}{factorialGapSegment_log_sum_le_channelLCM_add_choose}.}

\begin{proof}
Apply Lemma~\ref{long68:res:product-lcm} to $x_n=d_n$ for $N-k+1\le n\le N$.  Their
least common multiple divides $\Lc_N$.  By Lemma~\ref{long68:res:gap-gcd} each
pairwise gcd is smaller than $N^{\,j-i}$, and over a block of $k$ consecutive
indices
\[
 \sum_{i<j}(j-i)=\sum_{d=1}^{k-1}d(k-d)=\binom{k+1}{3}.
\]
Taking logarithms of the resulting divisibility gives \eqref{long68:eq:segment}.
\end{proof}

\begin{theorem}[common-denominator growth]\label{long68:res:lcm-growth}
\[
 \liminf_{N\to\infty}\frac{\log \Lc_N}{N^{3/2}\log N}
 \ \ge\ \frac{2\sqrt2}{3}.
\]
\end{theorem}
\leannote{Lean: \lproof{ErdosProblems/Erdos68/PaperCompleteLiminf.lean}{42}{common_denominator_growth_liminf}.}

\begin{proof}
A terminal block keeps every factorial near $N!$, but increasing its
length also increases the total pairwise-gcd loss.  The gain is of order
$kN\log N$ and the loss of order $k^3\log N$, so the useful balance is
$k$ of order $\sqrt N$.  Fix $\alpha>0$ and take
$k=\lfloor\alpha\sqrt N\rfloor$, which satisfies
$2\le k\le N-1$ for all large $N$.  Put $u=N-k+1$.

For the left-hand side of \eqref{long68:eq:segment}, use $n!-1\ge n!/2$ and
$\log n!\ge n\log n-n$.  Since $n\mapsto n\log n-n$ increases,
\[
 \sum_{n=u}^{N}\log(n!-1)\ \ge\ k\,(u\log u-u)-k\log2 .
\]
Here $u\ge N-\alpha\sqrt N$ and $k=\alpha\sqrt N+O(1)$, so
$k\,u\log u=\alpha N^{3/2}\log N+O(N\log N)$, while $k\,u=O(N^{3/2})$ and
$k\log2=O(\sqrt N)$.  Hence the left-hand side is at least
$\alpha N^{3/2}\log N+o(N^{3/2}\log N)$.

For the right-hand side, $\binom{k+1}{3}\le(k+1)^3/6$, so
$\binom{k+1}{3}\log N=\tfrac{\alpha^3}{6}N^{3/2}\log N+o(N^{3/2}\log N)$.
Therefore
\[
 \liminf_{N\to\infty}\frac{\log \Lc_N}{N^{3/2}\log N}
 \ \ge\ \alpha-\frac{\alpha^3}{6}.
\]
The right-hand side is maximised at $\alpha=\sqrt2$, with value
$\sqrt2-\tfrac{2\sqrt2}{6}=\tfrac{2\sqrt2}{3}$.
\end{proof}

The displayed proof of the lcm bound uses only
Lemmas~\ref{long68:res:product-lcm} and~\ref{long68:res:gap-gcd}.
Its formal counterpart is listed in the sources section; the status of that
entry is not a blanket assertion about all linked modules.
Theorem~12 of Garaev, Luca and Shparlinski
\cite[arXiv v1, p.~16]{garaev-luca-shparlinski} supplies a different input:
a uniform $O(N^{2/3})$ bound for the number of solutions of
$n!\equiv a\pmod p$, for fixed $a\not\equiv0\pmod p$, in an interval of
length $N$ contained in $1\le n<p$.  The weaker lcm deduction
from it is retained in \S\ref{long68:sec:ext-superseded}, not used here.
No exhaustive priority claim is made for the lcm estimate;
its precise elementary derivation and the inherited subtraction are stated
so that the claim can be assessed without relying on a search conclusion.

\begin{remark}[polynomial shifts]
Fix $P\in\Z[X]\smallsetminus\{0\}$.  By Lemma~3 of Luca and
Shparlinski~\cite{luca-shparlinski}, in the form given by
Lai~\cite[Lemma~2.1]{lai}, there is $n_0\ge2$, depending only on $P$, such
that $n!+P(n)>1$ for all $n\ge n_0$ and $P(n)(n+1)\cdots(n+h)\ne P(n+h)$ for
all $n\ge n_0$ and $h\ge1$. The nonvanishing has a short direct proof.
Beyond a cutoff, $P$ has constant nonzero sign and
$P(n+1)/P(n)\to1$, so $0<P(n+1)/P(n)<n+1$ for every sufficiently
large $n$. Multiplying these inequalities gives
\[
 0<\frac{P(n+h)}{P(n)}<\prod_{j=1}^{h}(n+j)\qquad(h\ge1).
\]
Thus the required subtraction is nonzero for every gap $h$, not merely
for each fixed $h$ after a separate cutoff. Increasing $n_0$ also makes
$n!+P(n)>1$. The proof of Theorem~\ref{long68:res:lcm-growth} then gives
\[
 \liminf_{N\to\infty}
 \frac{\log\lcm\{n!+P(n):n_0\le n\le N\}}{N^{3/2}\log N}
 \ \ge\ \frac{2\sqrt2}{3}.
\]
For $n_0\le i<j\le N$, multiplying $i!+P(i)$ by $j!/i!$ and subtracting
$j!+P(j)$ gives
\[
 \gcd\bigl(i!+P(i),\,j!+P(j)\bigr)\ \Big|\ \frac{j!}{i!}P(i)-P(j),
\]
which replaces Lemma~\ref{long68:res:gap-gcd}.  The right-hand side is
nonzero by the choice of $n_0$, and its absolute value is at most
$C_PN^{\,j-i+\deg P}$ for a constant $C_P$.  Over a block of $k$ indices the
pairwise loss therefore grows by $O_P(k^2\log N)$, which is $O_P(N\log N)$
when $k=\lfloor\alpha\sqrt N\rfloor$.  Since $n!+P(n)\ge n!/2$ for large $n$,
the lower estimate for the sum of logarithms is unchanged.  The cutoff is
needed, since for $P=-2$ the factor at $n=2$ vanishes.  The hypothesis
$P\ne0$ is needed as well, since for $P=0$ the least common multiple is $N!$
whatever the cutoff.
\end{remark}

\paragraph{A non-polynomial comparison.}
For $n!+2^n-1$, Luca and Shparlinski
\cite[Lemmas~2.1--2.3, pp.~860--862]{luca-shparlinski-exp} eliminate the
exponential term using three hits; multiplicative order enters the counting.
Their valuation-layer identity \cite[(3.2), p.~863]{luca-shparlinski-exp}
records repeated prime-power divisibility.  These are related tools, not
inputs to Theorem~\ref{long68:res:lcm-growth}, and no extension to general
non-polynomial perturbations is asserted.

Within the terminal-block estimate just proved, maximising
$\alpha-\alpha^3/6$ gives $2\sqrt2/3$.  This is an optimisation of that
one-parameter bound, not an optimality theorem for all block selections or
all gcd arguments.  The discarded factorial term has size $N^{3/2}$,
so a finite normalised ratio need not be close to the limiting lower bound.

Theorem~\ref{long68:res:lcm-growth} concerns $\Lc_N$, not $\den(H_N)$.
Theorem~\ref{long68:res:prime-pole} and the cancellations at $139$ and $2593$
show why these cannot be interchanged: a prime present in the common
denominator may be absent after reduction. The two examples establish this
possibility, not an asymptotic estimate for cancellation.

Even two consecutive denominators obstruct clearing by the full lcm. For
$N\ge3$, the gcd of $(N-1)!-1$ and $N!-1$ divides $N-1$, whereas
$(N-1)!-1$ is coprime to $N-1$. Thus those denominators are coprime, and
\[
\Lc_N(\Ser-H_N)>
\frac{((N-1)!-1)(N!-1)}{(N+1)!-1}\longrightarrow\infty.
\]
This argument does not use the stronger asymptotic growth theorem.

The direction of denominator control matters.  Write $H_N=\widehat A_N/Q_N$ in lowest
terms, reserving $A_N$ above for the numerator over the full common denominator.  Under the hypothetical identity $\Ser=a/q$, with $a\in\Z$ and
$q\ge1$, positivity of the tail gives
\[
 qQ_N(\Ser-H_N)=aQ_N-q\widehat A_N\in\Z_{>0}.
\]
A direct clearing contradiction therefore requires an upper estimate that
makes $qQ_N(\Ser-H_N)<1$ for some $N$; equivalently, it needs $Q_N$ to be
small relative to the reciprocal tail.  A lower bound on $Q_N$ points in the
opposite direction.  Indeed, the rational limit $1$ and the approximants
$1-1/Q_N$ have arbitrarily large reduced denominators without any
irrationality.  The inequality $d_N(\Ser-H_N)<1$ alone would not repair the
argument, because $d_N$ generally does not clear the earlier summands of
$H_N$.  The prime-power analysis describes which factors survive reduction,
but it does not yet supply the required upper control on $Q_N$.

For linear forms in $p$-adic zeta values, with $p\ge5$ a fixed prime and $n$
the index of the form, Lai, Lupu and Sprang quantify a saving of this kind.
They bound the denominators of the coefficients by powers of
$\lcm(1,\ldots,n)$, show that the bound may be divided by a product $\Phi_n$
of prime powers, and compute the growth rate of
$\Phi_n$~\cite[Lemmas~5.3--5.7 and~7.3]{lai-lupu-sprang}.  That saving enters
the inequality~\cite[(8.1)]{lai-lupu-sprang} under which the rescaled forms,
which have integer coefficients, meet the irrationality criterion
\cite[Lemma~2.1, p.~3]{lai-lupu-sprang} that they quote from Lai
\cite[Lemma~2.1, p.~4]{lai-2adic}. Along one unbounded subsequence, the forms
must be nonzero and their $p$-adic absolute values, multiplied by the
largest ordinary absolute value of an integer coefficient, must tend to
zero. The coefficient bound is archimedean; the bound on the value is
$p$-adic. The conclusion is that at least one of the finitely many
$p$-adic numbers in the forms is irrational, not that each is irrational.
The corresponding saving for $H_N$ is the ratio $\Lc_N/Q_N$.

This ratio is an integer.  Merely naming it supplies neither its growth
nor a small linear form; the analogy is about what an eventual estimate
would have to accomplish.
Theorem~\ref{long68:res:prime-pole} decides, for each prime $p\mid\Lc_N$,
whether $p$ divides this ratio; no asymptotic lower bound for the ratio is
established here.  This paragraph compares the two methods; no $p$-adic
theorem is applied to $\Ser$.

\section{Rationality and the next integer above a scaled partial sum}\label{long68:sec:carry}

For each scaled partial sum, take the least integer strictly above it.
This differs from the ordinary ceiling when the scaled sum is an integer.
Write
\[
 Z_m=\lfloor m!H_m\rfloor+1,\qquad
 \Delta_m=Z_{m-1}-(m-1)!H_{m-1}\quad(m\ge3).
\]
Thus $\Delta_m$ is the distance from the preceding scaled partial sum
to the next integer, and $0<\Delta_m\le1$. Define an integer $b_m$ by
\begin{equation}\label{long68:eq:carry-recurrence}
 Z_m=mZ_{m-1}+1-b_m .
\end{equation}


Call $b_m=1$ a \emph{unit carry}.  Put
$E_m=m!(\Ser-H_m)$ and $\varepsilon_m=1/(m!-1)$.  Since
$d_{n+1}>(n+1)d_n$, a geometric majorant gives the tail estimate
\begin{equation}\label{long68:eq:tail-bound}
 0<E_m<\frac{2\,m!}{(m+1)!-1}<\frac2m\le1\qquad(m\ge2).
\end{equation}
Compare the factorial scaling in Han\v{c}l and Tijdeman's tail-integrality
lemma for factorial series with integer coefficients
\cite[Lemma~2.1 and the following remark,
p.~385]{hancl-tijdeman}. Their scaled partial sums are integers; ours
need not be, since already $3!H_3=36/5$. The next proof instead uses the
short positive tail to identify the least integer above a scaled prefix
under a rationality assumption. The carry-defect expansion in
\S\ref{long68:sec:ext-literature} gives a separate, direct application of
factorial-series rationality criteria.

\label{long68:res:lead-carry-equivalence}\phantomsection
\label{long68:res:strict-successor-complete-characterization}\phantomsection
\begin{theorem}[an exact criterion from successive partial sums]\label{long68:res:carry-equivalence}
For $m\ge3$,
\begin{equation}\label{long68:eq:unit-window}
 b_m=1
 \iff m\mid Z_m
 \iff 1+\varepsilon_m<m\Delta_m\le2+\varepsilon_m .
\end{equation}


Moreover
\begin{align}
 \Ser\in\Q
 &\iff b_m=1\ \hbox{for all sufficiently large }m,
 \label{long68:eq:carry-rationality}\\
 \Ser\notin\Q
 &\iff \forall B\ \exists m>B:\ m\nmid Z_m .
 \label{long68:eq:strict-misses}
\end{align}
If $\Ser=a/q$ with $a\in\Z$, $q\ge1$ and $b_m\ne1$, then $q\nmid(m-1)!$ and
$q\ge m$.
\end{theorem}
\leannote{Lean: \lproof{ErdosProblems/Erdos68/PaperCompleteExisting.lean}{77}{strict_successor_characterisation}.}

\begin{proof}
From $m!H_m=mZ_{m-1}-m\Delta_m+1+\varepsilon_m$ one gets
$b_m=\lceil m\Delta_m-1-\varepsilon_m\rceil$ with $-1\le b_m\le m-1$, which
gives both equivalences in \eqref{long68:eq:unit-window}, endpoints included.

Suppose $\Ser=a/q$ and $q\mid(m-1)!$.  For $j=m-1$ and $j=m$ the number
$j!\,\Ser$ is an integer, and $j!H_j=j!\,\Ser-E_j$ with $0<E_j<1$ by
\eqref{long68:eq:tail-bound}, so $Z_j=j!\,\Ser$.  Substituting into
\eqref{long68:eq:carry-recurrence} forces $b_m=1$.  Every fixed $q$ divides
$(m-1)!$ eventually, which gives one direction of
\eqref{long68:eq:carry-rationality} and, by contraposition, the final assertion;
$q<m$ would imply $q\mid(m-1)!$.  Conversely, an eventual unit-carry tail
makes $Z_m/m!$ eventually constant, and $0<Z_m/m!-H_m\le1/m!$ with
$H_m\to\Ser$ identifies that constant as $\Ser$, which is then rational.
Negating \eqref{long68:eq:carry-rationality} gives \eqref{long68:eq:strict-misses}.
\end{proof}

The conclusion $q\nmid(m-1)!$ is about prime-power multiplicities, not
just prime factors. It permits all prime factors of $q$ to be at most
$m-1$ if one occurs to a larger exponent than in $(m-1)!$.

\subsection{Why one proposed window argument is circular}
\label{long68:sec:adjacent-unit-no-go}

A proposed argument starts by assuming $b_m=b_{m+1}=1$, clears the
denominators in the resulting inequalities, and tries to contradict the
size of a prime-power divisor. The following calculation shows why the
resulting size bound alone cannot provide that contradiction.

Write $\Delta_k=P^\Delta_k/Q^\Delta_k$ in lowest terms, with $P^\Delta_k,Q^\Delta_k>0$.
In one recurrence step the unreduced denominator is $Q^\Delta_k(k!-1)$;
let $G_k=Q^\Delta_k(k!-1)/Q^\Delta_{k+1}$ be the positive integer cancelled in reduction.
For $m\ge3$, put
\[
 \begin{aligned}
 D_m&=Q^\Delta_m(m!-1)((m+1)!-1),\\
 \Omega_m&=D_m\left(m(m+1)\Delta_m-m-2
       -\frac{m+1}{m!-1}-\frac1{(m+1)!-1}\right).
 \end{aligned}
\]
Both are integers. The two unit carries are equivalent to
$0<\Omega_m\le D_m$. Independently of the carries, cancellation in two
successive steps gives
\[
 D_m=Q^\Delta_{m+2}G_{m+1}G_m=Q^\Delta_m(m!-1)\bigl((m+1)!-1\bigr),
\]
and under the pair assumption the offset factors to match,
$\Omega_m=P^\Delta_{m+2}G_{m+1}G_m$.  Dividing by $G_{m+1}G_m>0$ reduces the proposed bound to
\[
 0<P^\Delta_{m+2}\le Q^\Delta_{m+2},
\]
which every reduced positive gap in $(0,1]$ already satisfies.  Thus the resulting gap bound is not an additional restriction on a
pair of unit carries. This calculation does not exclude an independent
arithmetic obstruction involving the transition numerators or their
valuations. Such an obstruction would need information not already
implied by the pair assumption; it need not take the form of a bound on
the preceding gap $\Delta_m$.

\subsection{Factorial digits of \texorpdfstring{$S-e+2$}{S-e+2}}

Let $C=\sum_{n\ge2}\bigl(n!(n!-1)\bigr)^{-1}$.  The identity
$1/(n!-1)=1/n!+1/(n!(n!-1))$ and absolute convergence give
\begin{equation}\label{long68:eq:companion-decomposition}
 \Ser=C+e-2 .
\end{equation}
The canonical factorial digits of a real $x$ are
$a_m(x)=\lfloor m!x\rfloor-m\lfloor(m-1)!x\rfloor\in\{0,\ldots,m-1\}$
for $m\ge2$. At a rational endpoint, the floor convention selects the
terminating representation. For example, $1/2=1/2!$ also equals
$\sum_{m\ge3}(m-1)/m!$, but its canonical digits are $a_2=1$ and
$a_m=0$ for $m\ge3$, not an eventually maximal tail. With this convention,
$x$ is rational exactly when its digits vanish from some index on. The criterion
goes back to Cantor \cite{cantor1869}; Galambos treats the rationality of
Cantor series in \cite[Ch.~II, \S2.1, pp.~21--22]{galambos1976}, and Koepf
and Schmersau prove the irrationality direction for nonterminating factorial
expansions whose digits are not eventually maximal
\cite[Example~3.2, p.~121]{koepf-schmersau}.

\label{long68:res:lead-companion-orbit}\phantomsection
\begin{proposition}[rationality and factorial residues]\label{long68:res:companion-orbit}
\[
 \Ser\in\Q
 \quad\Longleftrightarrow\quad
 \lfloor m!C\rfloor\equiv-2\pmod m
 \quad\hbox{for all sufficiently large }m.
\]
\end{proposition}
\leannote{Lean: \lproof{ErdosProblems/Erdos68/PaperCompleteExisting.lean}{38}{companion_orbit}.}

\begin{proof}
Put $J_m=\sum_{k=0}^{m}m!/k!$.  Then $J_m$ is an integer, every summand with
$k<m$ is divisible by $m$ and the summand at $k=m$ is $1$, so
$J_m\equiv1\pmod m$; and $0<m!\,e-J_m<1$ for $m\ge2$.

Suppose $\Ser$ is rational.  For all large $m$ both $(m-1)!\,\Ser$ and
$m!\,\Ser$ are integers, so $m\mid m!\,\Ser$.  Multiplying
\eqref{long68:eq:companion-decomposition} by $m!$ and using
$\lfloor -m!\,e\rfloor=-J_m-1$ gives
$\lfloor m!C\rfloor=m!\,\Ser+2\,m!-J_m-1\equiv-2\pmod m$.

Conversely, assume the congruence from some index on.  Because
$0\le a_m(C)<m$, for $m\ge3$ the congruence is equivalent to $a_m(C)=m-2$.
Choose $N$ beyond the exceptional indices.  The canonical expansion gives
\[
 C=\frac{\lfloor N!C\rfloor}{N!}+\sum_{m>N}\frac{m-2}{m!},
\]
and the telescope $\sum_{m>N}(m-1)/m!=1/N!$ turns this into
\[
 \Ser=C+e-2
 =\frac{\lfloor N!C\rfloor+1}{N!}+\sum_{m=2}^{N}\frac1{m!},
\]
which is rational.
\end{proof}

The condition is eventual equality, not equality at many computed
indices. The proof works for any real $C$: it holds exactly when
$C+e-2$ is rational. For example, $C=2-e$ gives a rational sum, whereas
$C=0$ gives the irrational number $e-2$. The issue here is to decide the
condition for the particular positive series defining $C$.

\subsection{Escape from a smaller interval}

Write $\theta_m=\{m!\,\Ser\}$. We ask whether $\theta_{m-1}$ lies outside
$[0,E_m/m)$, an interval of width less than $2/m^2$.

\begin{proposition}[lower-interval criterion]\label{long68:res:lower-escape}
\begin{equation}\label{long68:eq:lower-escape}
 \Ser\notin\Q
 \quad\Longleftrightarrow\quad
 \forall B\ \exists m>B:\ E_m\le m\theta_{m-1}.
\end{equation}
For $m\ge3$ the finite condition
\begin{equation}\label{long68:eq:finite-escape}
 m\Delta_m\le1+\varepsilon_m
 \quad\hbox{or}\quad
 1+\varepsilon_m+\frac2m\le m\Delta_m
\end{equation}
implies the escape inequality in \eqref{long68:eq:lower-escape}.  Cofinally many
instances of \eqref{long68:eq:finite-escape} therefore imply $\Ser\notin\Q$.
\end{proposition}
\leannote{Lean: \lproof{ErdosProblems/Erdos68/PaperCompleteExisting.lean}{102}{lower_interval_criterion}.}

\begin{proof}
If $\Ser$ is rational, the remainders $\theta_{m-1}$ vanish eventually while
$E_m>0$, so the escape inequality fails eventually.  Conversely, suppose
$m\theta_{m-1}<E_m$ for every sufficiently large $m$.  By
\eqref{long68:eq:tail-bound} this gives $m\theta_{m-1}<1$, so $a_m(\Ser)=0$ and
$\theta_m=m\theta_{m-1}$. Fix $N$ after this recurrence begins.
Then $\theta_{N+k}=((N+k)!/N!)\theta_N<1$ for every $k\ge0$, forcing
$\theta_N=0$. Hence $N!\Ser$ is an integer and $\Ser$ is rational.

For the finite implication, suppose $m\theta_{m-1}<E_m$.  The tail recurrence
is $mE_{m-1}=1+\varepsilon_m+E_m$, so $E_{m-1}>E_m/m>\theta_{m-1}\ge0$, and
with $0<E_{m-1}<1$ the strict-successor identity gives
$\Delta_m=E_{m-1}-\theta_{m-1}$.  Hence
\[
 1+\varepsilon_m<m\Delta_m=1+\varepsilon_m+E_m-m\theta_{m-1}
 \le1+\varepsilon_m+E_m<1+\varepsilon_m+\frac2m,
\]
which contradicts \eqref{long68:eq:finite-escape}.
\end{proof}

Escape at a single index does not give a non-unit carry there.

The finite test excludes only the open interval
$1+\varepsilon_m<m\Delta_m<1+\varepsilon_m+2/m$; it accepts either endpoint
and all values outside. The width $2/m$ refers to $m\Delta_m$; the interval
for $\Delta_m$ has width $2/m^2$. At one index this is a sufficient test
for escape, obtained by replacing $E_m$ with the upper bound $2/m$.

At arbitrarily large integer indices, however, the finite test itself
characterises irrationality:
\[
 S\notin\Q\quad\Longleftrightarrow\quad
 \text{\eqref{long68:eq:finite-escape} holds for arbitrarily large integers }m.
\]
Only the forward implication remains to be shown. If $S$ is irrational,
Theorem~\ref{long68:res:carry-equivalence} gives $b_m\ne1$ at arbitrarily
large indices. By \eqref{long68:eq:unit-window}, each such index satisfies
$m\Delta_m\le1+\varepsilon_m$ or $m\Delta_m>2+\varepsilon_m$.
Since $2/m<1$ for $m\ge3$, either alternative implies
\eqref{long68:eq:finite-escape}. The converse is Proposition~\ref{long68:res:lower-escape}.
This is an ordinary consequence of the two proved criteria, not a proof
that either occurs infinitely often for $S$. The equivalence is over all
integer indices; restricting to primes is only a sufficient condition here.

For comparison at a single index, the exact classification is
\begin{equation}\label{long68:eq:escape-classification}
 E_m\le m\theta_{m-1}
 \quad\Longleftrightarrow\quad
 b_m\ne1\ \hbox{ or }\ a_m(\Ser)=m-1
 \qquad(m\ge3).
\end{equation}

To verify the classification, use $0<E_j<1$ to write
\[
 Z_j=\lfloor j!S\rfloor+\mathbf1_{\{\theta_j\ge E_j\}}.
\]
Substitution in the carry recurrence gives
\[
 b_m=1-a_m(S)+m\mathbf1_{\{\theta_{m-1}\ge E_{m-1}\}}
                 -\mathbf1_{\{\theta_m\ge E_m\}}.
\]
Since $0\le a_m(S)\le m-1$, a unit carry has exactly two possibilities:
both indicators are zero and $a_m(S)=0$, or both are one and
$a_m(S)=m-1$. In the first case $m\theta_{m-1}=\theta_m<E_m$;
in the second, $m\theta_{m-1}\ge m-1>E_m$.
Conversely, $m\theta_{m-1}<E_m<1$ forces $a_m(S)=0$ and both
indicators to be zero, using
$mE_{m-1}=1+\varepsilon_m+E_m$. It therefore forces a unit carry with
zero digit. This proves \eqref{long68:eq:escape-classification}, including
the equality case in its left-hand inequality.

Thus escape can occur even when $b_m=1$, provided the digit is $m-1$.
This happens at $m=52$: exact
rational arithmetic gives $b_{52}=1$ together with
$52\,\Delta_{52}>1+\varepsilon_{52}+2/52$, the margin being about
$0.5689674908$.  This example separates the two events at a single index: the finite test
can hold at a unit carry. Their occurrence at arbitrarily large indices
nevertheless gives equivalent rationality criteria. Since $52$ is
composite, the example supplies no prime-index instance.

\subsection{The series with denominators \texorpdfstring{$n!+t$}{n!+t}}

The same digit argument applies to the other shifts mentioned by
Erd\H{o}s.  For an integer $t\ge-1$ put
$\Ser_t=\sum_{n\ge2}1/(n!+t)$ and $C_t=\sum_{n\ge2}1/\bigl(n!(n!+t)\bigr)$,
so that $\Ser_t=-tC_t+(e-2)$.

The restriction $t\ge-1$ keeps every denominator positive for $n\ge2$;
$t=-2$ is excluded because its first denominator is zero. All series in
this identity converge absolutely. The familiar case $t=0$ is $e-2$;
this example will make the residue condition explicit.

\begin{theorem}[a criterion for the shifts $t\ge-1$]\label{long68:res:shift-family}
For every integer $t\ge-1$, the series $\Ser_t$ is rational exactly when
\[
 \bigl\lceil t\,m!\,C_t\bigr\rceil\equiv2\pmod m
\]
for all sufficiently large $m$, and irrational exactly when that residue is
missed cofinally.  The member $t=-1$ is $\Ser$, and the member $t=0$ is
$e-2$, for which the scaled correction is always $0$ and hence misses the residue class
at every $m\ge3$, proving the irrationality of $e$.
\end{theorem}
\leannote{Lean: \lproof{ErdosProblems/Erdos68/PaperCompleteExisting.lean}{123}{uniform_family_boundary}, \lproof{ErdosProblems/Erdos68/PaperCompleteExisting.lean}{134}{uniform_family_members}.}

\begin{proof}
Put $Y=-tC_t$, so that $\Ser_t=Y+e-2$.  The proof of
Proposition~\ref{long68:res:companion-orbit} applies to any real $Y$:
rationality of $Y+e-2$ forces
$\lfloor m!Y\rfloor\equiv-2\pmod m$ eventually.  Conversely, that congruence
forces the canonical digits of $Y$ to equal $m-2$ eventually; adding the
factorial series of $e-2$ leaves an eventually telescoping tail, hence a
rational sum.  Finally
$\lceil t m!C_t\rceil=-\lfloor m!Y\rfloor$ translates the congruence.
Negation gives the cofinal statement.  For $t=0$, the left side is zero,
which is not congruent to $2$ for $m\ge3$.
\end{proof}

The $t=0$ specialisation recovers the classical irrationality of $e$.
For $t\ne0$, no estimate proving failure of the displayed congruence at
arbitrarily large indices is supplied here.
This is a statement about the scope of the present argument, not an assertion
that every other member has the same current literature status.

\section{A sufficient comparison between the tail and an integer gap}\label{long68:sec:residue}

We choose to clear every prime-power level shared by two summand
denominators. Each remaining prime then occurs in just one scaled
summand denominator, so its maximal exponent cannot cancel. This is a
sufficient construction, not a necessary size for a clearing factor:
cancellation in the partial sum can make a smaller scale suffice.
The remaining denominator will determine the gap to the next integer. For an integer $p\ge3$, not necessarily prime, put
$I_p=\{2,\ldots,2p-1\}$ and $F_p=(p-1)!$, and define
\begin{align*}
 D_p&=\lcm_{\substack{i,j\in I_p\\ i<j}}\gcd(d_i,d_j),
 &C_p&=\lcm(F_p,D_p),\\
 \Lc^{\mathrm{blk}}_p&=\lcm(F_p,d_2,\ldots,d_{2p-1}),
 &R_p&=\Lc^{\mathrm{blk}}_p/C_p .
\end{align*}
The integer $D_p$ contains the prime powers shared by at least two
of the denominators. The integer $C_p$ also contains the factorial base
$F_p$, and $R_p$ is the quotient left in the full common denominator.
The factors $C_p$ and $R_p$ need not be coprime: a prime absent from $F_p$
with largest exponent $5$ and second-largest exponent $2$ occurs to powers
$2$ and $3$ in the two factors. This is an illustration of the valuation
calculation, not a claimed pattern in a particular factorial block.
It will be useful to retain the ratio $\widetilde C_p=C_p/F_p$, so that
\begin{equation}\label{long68:eq:core-normalisation}
 \Lc^{\mathrm{blk}}_p=C_pR_p=F_p\widetilde C_pR_p .
\end{equation}


Including $F_p$ in both $C_p$ and $\Lc_p^{\mathrm{blk}}$ ensures that
a fixed rational denominator divides $C_p$ for all large $p$.
The ratio $\widetilde C_p$ removes precisely that factorial factor when
we estimate the shared part. Put
\[
 T_p=\sum_{i\in I_p}\frac{\Lc^{\mathrm{blk}}_p}{d_i},\qquad
 \rho_p=(-T_p)\bmod R_p,\qquad
 K_p=2p^2(2p-1)! ,
\]
with $\rho_p$ the least nonnegative representative.  All of these depend only
on a finite prefix.

For each prime, $C_p$ contains the larger of its exponent in $F_p$ and
its second-largest exponent among the denominators.  A prime dividing $R_p$ therefore has a unique
denominator whose exponent is maximal and exceeds the exponent in $F_p$, so the argument of
Theorem~\ref{long68:res:prime-pole} applies to it and gives $\gcd(T_p,R_p)=1$.
Thus $R_p=\den(C_pH_{2p-1})$, the reduced denominator of the scaled
prefix, not generally of $H_{2p-1}$. When $R_p>1$, this ratio is
nonintegral and
\begin{equation}\label{long68:eq:gap-normalisation}
 \frac{\rho_p}{R_p}
 =\bigl\lceil C_pH_{2p-1}\bigr\rceil-C_pH_{2p-1}\in(0,1),
\end{equation}
the distance to the next integer. When $R_p=1$, $\rho_p=0$ but the
scaled partial sum is integral, so the least strictly larger integer is at
distance $1$. The theorem below uses only $R_p>1$.

At $p=3$, the denominators $1,5,23,119$ are pairwise coprime. Thus $D_3=1$,
$C_3=2$, $R_3=13685$, and
\[
C_3H_5=\frac{34264}{13685},\qquad
\frac{\rho_3}{R_3}=3-\frac{34264}{13685}=\frac{6791}{13685}.
\]
The tail estimate in the proof below gives
$C_3(\Ser-H_5)<7/1080<6791/13685$. This verifies the comparison for one
actual prefix. The theorem needs such prefixes at arbitrarily large $p$,
so that $F_p$ can absorb any fixed rational denominator.

\begin{theorem}[a sufficient tail inequality]\label{long68:res:global-residue}
Suppose that for every $B$ there is a natural parameter $p\ge3$ with
$p>B$, $R_p>1$, and
\begin{equation}\label{long68:eq:global-scale}
 (2p+1)\Lc^{\mathrm{blk}}_p<K_p\rho_p .
\end{equation}
Then $\Ser$ is irrational.
\end{theorem}
\leannote{Lean: \lproof{ErdosProblems/Erdos68/PaperCompleteExisting.lean}{169}{global_complementary_criterion_prime}.}

\begin{proof}
Suppose $\Ser=a/q$ with $q\ge1$, and choose a parameter in the hypothesis
with $p>q$. Since $q\mid F_p\mid C_p$, the number $C_p\Ser$ is an integer
strictly above $C_pH_{2p-1}$. The gap identity
\eqref{long68:eq:gap-normalisation} gives
\[
 C_p(\Ser-H_{2p-1})\ge\frac{\rho_p}{R_p}.
\]
On the other hand, $1/(n!-1)<2/n!$ for $n\ge2p$, and comparison with a
geometric series of ratio $1/(2p+1)$ gives
\[
 \Ser-H_{2p-1}<\frac{2}{(2p)!}\frac{2p+1}{2p}=\frac{2p+1}{K_p}.
\]
Multiplication by $C_p$ and \eqref{long68:eq:global-scale}, with
$\Lc_p^{\mathrm{blk}}=C_pR_p$, now give
$C_p(\Ser-H_{2p-1})<\rho_p/R_p$, a contradiction.
\end{proof}

\noindent The formal statement for natural-number parameters is
\href{https://github.com/wcook04/plectis-erdos/blob/f36a98bf3d3e6f65f1074e3b800e3293d5b8a51a/ErdosProblems/Erdos68/PaperCompleteExisting.lean\#L156-L168}{comparison of the tail with the next integer}.

Cancelling $R_p$ in \eqref{long68:eq:global-scale} through
\eqref{long68:eq:core-normalisation} displays the quantitative issue directly:
\begin{equation}\label{long68:eq:core-gap-scale}
 (2p+1)C_p<K_p\frac{\rho_p}{R_p}.
\end{equation}
The factor $R_p$ has cancelled. Thus new prime divisors, even large ones,
do not by themselves bound $C_p$ or keep $\rho_p/R_p$ away from zero.
Proposition~\ref{long68:res:wilson-cofinality} supplies the primes but not
these two estimates.

Under $R_p>1$, we have $\rho_p>0$, and taking logarithms gives
\[
 \log\widetilde C_p-\log\frac{\rho_p}{R_p}
 <\log\frac{2p^2(2p-1)!}{(2p+1)(p-1)!}
 =p\log p+(2\log2-1)p+O(\log p).
\]
Both terms on the left must be controlled at the same $p$. The elementary
bound $\rho_p\ge1$ gives only $\rho_p/R_p\ge1/R_p$, which may be too small.
For a concrete comparison, a gap of at least $1/2$ would reduce the sufficient
test to $(2p+1)C_p<K_p/2$. At the opposite extreme, the gap $1/R_p$ requires
the stronger bound $(2p+1)C_p<K_p/R_p$. Neither type of favourable behaviour
is proved here for an unbounded family of these particular prefixes.

\label{long68:res:lead-moving-factor-split}\phantomsection
\label{long68:res:moving-factor-scale-split}\phantomsection
A more restrictive test also uses a prime $q$ that first divides $m!-1$
at an index $m\ge4$. Choose $p=\lfloor m/2\rfloor+1$, which need not be
prime. This choice puts $m$ in the block and gives $q\mid R_p$. Indeed,
$q>m\ge p$, so $q\nmid F_p$, and no earlier denominator contains $q$.
The only possible later index in the block is $m+1$; when it occurs,
$(m+1)!-1\equiv m\not\equiv0\pmod q$. Thus $q$ occurs in just one block
denominator and is absent from $C_p$.

The additional comparison uses the integer $R_p/q$. Requiring the tail
bound to be smaller than both $\rho_p$ and $R_p/q$ gives
$\min(\rho_p,R_p/q)$ and is equivalent to the following two inequalities:
\begin{equation}\label{long68:eq:factor-split}
 (2p+1)\Lc^{\mathrm{blk}}_p<K_p\min(\rho_p,R_p/q)
 \quad\Longleftrightarrow\quad
 \begin{cases}
 (2p+1)\Lc^{\mathrm{blk}}_p<K_p\rho_p,\\
 (2p+1)C_pq<K_p .
 \end{cases}
\end{equation}


The equivalence follows by testing each entry of the minimum and using
\eqref{long68:eq:core-normalisation}. Both inequalities must hold for the
same arbitrarily large parameters. The first alone is already sufficient
by Theorem~\ref{long68:res:global-residue}; the second is an extra restriction,
not an equivalent formulation of irrationality.

\label{long68:res:split-factor-normalized-collision}\phantomsection
\label{long68:bdry:fixed-owner-absorption}\phantomsection
A fixed small index cannot supply new factors indefinitely. If $n$ is
fixed and $p>n!-1$, then $n!-1$ divides $(p-1)!$, so none of its prime
powers remains outside the factorial base. For example, the factor $719$
of $6!-1$ is useful only over a finite range of $p$. Any argument using
such factors at arbitrarily large $p$ must let their indices increase.

\section{What can be achieved by cancelling finitely many weighted sums}\label{long68:sec:channels}

We next try to remove the first few terms of a remainder by integer linear
combinations. The weights are chosen so that their difference from ordinary
factorial weights is divisible by $d!-1$. This preserves fractional parts
after division by $d!-1$, while allowing the contributions at selected
denominators to vanish.
For a finitely supported integer vector $c=(c_i)$ on indices $i\ge2$, put
\[
 M(c)=\sum_i c_i\,i!,\qquad
 W_{d,i}=\frac{i!}{(d!)^{\lfloor i/d\rfloor}},\qquad
 \chn dc=\sum_i c_iW_{d,i}\quad(d\ge2).
\]
We call $M(c)$ the moment; it and the weighted sums $V_d(c)$ may be
negative. The remainder is
\[
 \mathcal R(c)=\sum_{d\ge2}\frac{V_d(c)}{d!-1}.
\]
Beyond the support, $V_d(c)=M(c)$, so this series converges absolutely.
Cancelling its first weighted sums leaves a finite signed block followed
by $M(c)$ times the original tail. The congruence below explains why this
is still an integer linear form in $1$ and $S$.
The power $(d!)^{\lfloor i/d\rfloor}$ need not be the largest power
of $d!$ dividing $i!$: at $d=2$, $i=6$, the chosen exponent is $3$,
although $2^4\mid6!$. We use the floor exponent because it changes
between $i-1$ and $i$ exactly when $d\mid i$. Adjacent differences will therefore affect
only divisor-indexed weighted sums, which is the reason the elimination
in \S5.1 works.
Each $W_{d,i}$ is an integer: writing $i=kd+r$ with $0\le r<d$, the quotient
$i!/\bigl((d!)^kr!\bigr)$ is a multinomial coefficient.  Since
$i!=(d!)^{\lfloor i/d\rfloor}W_{d,i}$ and $d!\equiv1\pmod{d!-1}$,
\begin{equation}\label{long68:eq:channel-congruence}
 \chn dc\equiv M(c)\pmod{d!-1}.
\end{equation}

\begin{theorem}[divisibility of the difference]\label{long68:res:normalform}
For every finite integer support and every $d\ge2$ there is an integer $k$
with $\chn dc=M(c)+(d!-1)k$.
\end{theorem}
\leannote{Lean: \lproof{ErdosProblems/Erdos68/PaperCompleteSupportedBands.lean}{17}{supported_integral_normal_form}.}

\begin{proof}
By \eqref{long68:eq:channel-congruence}, $d!-1$ divides
$V_d(c)-M(c)$.  Thus $k=(V_d(c)-M(c))/(d!-1)$ is an integer and gives the
identity.  The modulus is positive also at $d=2$, when it equals one.
\end{proof}

For any $N\ge2$ at least as large as every supported index, the terms
with $d>N$ equal $M(c)/(d!-1)$. Therefore
\[
 \mathcal R(c)-M(c)\Ser
 =\sum_{d=2}^N\frac{V_d(c)-M(c)}{d!-1}\in\mathbb Z.
\]
The sum is finite and each term is integral by
\eqref{long68:eq:channel-congruence}. This argument also holds when
index $1$ is temporarily allowed in the basis calculation below.
In particular, a zero-moment vector has integral remainder, and fixed
moment fixes the fractional part. These facts do not depend on solving
the cancellation equations.

In the next theorem, the parameters $d,k$ and the support indices are integers.

\begin{theorem}[constant values of the floor in the weights]\label{long68:res:bandbreakpoint}
Let $d\ge2$ and $k\ge0$, and suppose every supported index $i$ satisfies
$kd\le i<(k+1)d$.  Then $M(c)=(d!)^k\chn dc$.  In particular, cancellation
on the interval $d\le i<2d$ forces $M(c)=0$; and if every supported index
is at least $d$ while $M(c)\ne0$ and $\chn dc=0$, then some supported index
is at least $2d$.
\end{theorem}
\leannote{Lean: \lproof{ErdosProblems/Erdos68/PaperCompleteSupportedBands.lean}{24}{supported_quotient_band}, \lproof{ErdosProblems/Erdos68/PaperCompleteSupportedBands.lean}{42}{supported_first_band_cancellation}, \lproof{ErdosProblems/Erdos68/PaperCompleteSupportedBands.lean}{51}{supported_breakpoint_escape}.}

\begin{proof}
On this interval the quotient $\lfloor i/d\rfloor$ is the constant $k$, so
$i!=(d!)^kW_{d,i}$ for every supported index and the factor $(d!)^k$ comes
out of the sum.  The two consequences follow by taking $k=1$ and by
contraposition.
\end{proof}

\noindent The hypothesis means that all supported indices lie in one
interval on which $\lfloor i/d\rfloor$ is constant. A single supported
index always has this property; indices on opposite sides of a multiple
of $d$ need not. The conclusion restricts the support of a solution with
$V_d(c)=0$ and $M(c)\ne0$. It neither constructs such a solution nor bounds
its infinite remainder.

By \eqref{long68:eq:channel-congruence}, a vanishing $d$-th weighted sum forces $(d!-1)\mid M(c)$,
and annihilating every weighted sum $2\le d\le D$ forces $\Lc_D\mid M(c)$, where
$\Lc_D=\lcm_{2\le d\le D}(d!-1)$ is the same quantity as in
\S\ref{long68:sec:lcm}. At fixed moment, the same congruence fixes the
fractional part of each quotient $V_d(c)/(d!-1)$ separately.

\subsection{All solutions and their remainders modulo integers}

Fix $D\ge2$. We will solve $V_2=\cdots=V_D=0$, impose the restriction
that the coefficient at index $1$ vanish, and then calculate the resulting
infinite remainder. The support restriction is an additional equation,
not a consequence of cancelling the weighted sums.

Temporarily allow finitely supported integer vectors on $n\ge1$, with
$M$ and $\chn d{\cdot}$ given by the same formulas; this introduces no term
$1/(1!-1)$ into $\Ser$. Let $e_n$ have coefficient $1$ at index $n$ and
zero elsewhere. Adjacent factorials suggest starting with
\[
 T_n=ne_{n-1}-e_n\quad(n\ge2),
\]
which has moment zero. The quotient $\lfloor i/d\rfloor$ changes between
$i=n-1$ and $i=n$ exactly when $d\mid n$, so
\[
 \chn d{T_n}=(d!-1)W_{d,n}\mathbf1_{d\mid n}.
\]
For example, $T_4=4e_3-e_4$ has $V_2(T_4)=6$ and $V_4(T_4)=23$, with
all other weighted sums zero. Since $U_2:=T_2=2e_1-e_2$ has only
$V_2(U_2)=1$ nonzero, subtracting $6U_2$ leaves only the fourth weighted
sum. This suggests removing the contribution of each proper divisor:
\[
 U_n=T_n-\sum_{\substack{d\mid n\\2\le d<n}}W_{d,n}U_d\quad(n\ge2).
\]
Induction gives $M(U_n)=0$ and
$\chn d{U_n}=(d!-1)\mathbf1_{d=n}$.
On each finite initial segment, the columns $e_1,T_2,T_3,\ldots$ form a
triangular integer matrix with diagonal entries $1,-1,-1,\ldots$.
Its determinant is $\pm1$, so it is unimodular. The change from $T_n$ to
$U_n$ is triangular over $\mathbb Z$ with diagonal entries $1$.
Thus $e_1,U_2,U_3,\ldots$ is a $\mathbb Z$-basis of the finitely supported
integer vectors. Applying $M$ and each $\chn d{\cdot}$ gives the unique expansion
\[
 c=M(c)\,e_1+\sum_{d\ge2}\frac{\chn dc-M(c)}{d!-1}\,U_d .
\]
When $V_2=\cdots=V_D=0$, the congruences above force $\Lc_D\mid M(c)$.
To obtain moment $\Lc_D$, define
\[
 K_D=\Lc_De_1-\sum_{d=2}^{D}\frac{\Lc_D}{d!-1}U_d.
\]
Then $M(K_D)=\Lc_D$ and $\chn d{K_D}=0$ for $2\le d\le D$.
Consequently every vector of the enlarged space whose weighted sums $2,\ldots,D$
vanish has a unique expression
\begin{equation}\label{long68:eq:low-channel-basis}
 c=tK_D+\sum_{n>D}z_nU_n,
 \qquad M(c)=tL_D,
\end{equation}
with finitely many nonzero integers $z_n$.  Let $a_D$ and $u_n$ denote the coefficients at index $1$ in $K_D$ and
$U_n$, respectively.  Requiring support on $n\ge2$ adds precisely the scalar
equation
\begin{equation}\label{long68:eq:low-channel-support}
 ta_D+\sum_{n>D}z_nu_n=0.
\end{equation}
The scalar coefficients satisfy $u_2=2$ and
$u_n=-\sum_{d\mid n,\ 2\le d<n}W_{d,n}u_d$ for $n>2$.
Every proper divisor of an odd $n$ is odd, so induction gives $u_n=0$
for odd $n>1$. At $n=2p$ with $p$ prime, only $d=2$ contributes a
nonzero term; hence $u_{2p}=-(2p)!/2^{p-1}$, also for $p=2$.
If $p$ is a prime with
$D/2<p\le D$ and $H=D(2p-1)$, then
\begin{equation}\label{long68:eq:finite-horizon}
 \gcd\{u_n:n>D\}=\gcd(u_{D+1},\ldots,u_H)>0,\qquad H<2D^2 .
\end{equation}
Indeed, the finite gcd $g$ divides $u_{2p}\ne0$ and hence $(2p)!$.  For $n>H$
and a divisor $d\le D$ of $n$ one has $n/d\ge2p$. The quotient
$W_{d,n}/(n/d)!$ counts partitions of $n$ points into $n/d$ unordered
blocks of size $d$, so $(n/d)!\mid W_{d,n}$ and hence $g\mid W_{d,n}$.  Strong induction in the
recurrence handles the remaining divisors $D<d<n$, so $g$ divides every $u_n$
with $n>D$.  Bertrand's postulate supplies $p$ for $D\ge3$, and $p=2$ serves
for $D=2$; finally $H\le D(2D-1)<2D^2$.  The short note proves the same
determination in its theorem ``a finite formula for the gcd''.

\begin{theorem}[the set of attainable moments]\label{long68:res:moment-ideal}
Fix $D\ge2$ and a prime $p$ with $D/2<p\le D$.  Put
\[
 H=D(2p-1),\qquad
 G_D=\gcd\{|u_n|:D<n\le H\},\qquad
 \mu_D=L_D\frac{G_D}{\gcd(G_D,a_D)}.
\]
The moments of finite integer vectors supported on $n\ge2$ and cancelling all
weighted sums $2,\ldots,D$ are exactly $\mu_D\Z$.  The integer $\mu_D$ is positive,
is independent of the eligible prime $p$, and is attained by a vector of
coefficients with gcd one, that is, by a primitive vector.
\end{theorem}
\leannote{Lean: \lproof{ErdosProblems/Erdos68/PaperCompleteMomentIdeal.lean}{242}{exact_moment_ideal_with_primitive_attainment}, \lproof{ErdosProblems/Erdos68/PaperCompleteMomentIdeal.lean}{321}{minimumMoment_independent_prime}.}

\begin{proof}
By \eqref{long68:eq:finite-horizon}, $G_D>0$ is the gcd of the whole tail
$\{u_n:n>D\}$, and $H<2D^2$.  Thus the finite integer
combinations in \eqref{long68:eq:low-channel-support} form $G_D\Z$, so that the
equation is soluble exactly when $G_D\mid ta_D$.  Dividing by
$\gcd(G_D,a_D)$ shows that $t$ is a multiple of
$G_D/\gcd(G_D,a_D)$, and \eqref{long68:eq:low-channel-basis} gives the displayed
moment ideal.  Bezout coefficients attain its positive generator.  If an
attaining vector had a nontrivial common coefficient divisor, division by that
divisor would produce a smaller positive attainable moment.  Hence its coefficients have gcd one. Since the set of attainable moments
does not depend on the chosen prime, neither does its least positive element.
\end{proof}

For a concrete instance, take $D=4$. One first needs
$\lcm(1,\ldots,n)\mid u_n$, not just a list of computed coefficients.
For $d\mid n$, Legendre's formula shows that the valuation of $W_{d,n}$
at a prime $r$ is at least the number of powers of $r$ in $(d,n]$:
each such power contributes at least $1$ to
$\sum_{j\ge1}(\lfloor n/r^j\rfloor-(n/d)\lfloor d/r^j\rfloor)$,
and the other terms are nonnegative. Hence
$\lcm(1,\ldots,n)\mid W_{d,n}\lcm(1,\ldots,d)$, and induction in the
recurrence for $u_n$ proves the divisibility. It follows that $60$ divides
every $u_n$ with $n>4$. Conversely,
$u_6=-180$, $u_8=-4200$ and $23u_6-u_8=60$, so $G_4=60$.
Since $L_4=115$ and $a_4=-55$, the least positive moment is therefore
$115\cdot60/\gcd(60,55)=1380$. The vector
$12K_4+253U_6-11U_8$ attains it: its coefficient at index $1$ is
$12(-55)+253(-180)-11(-4200)=0$. This example uses the gcd of the
entire tail, not only the observed coefficients.

If support is additionally restricted to $n\le6$, only $u_5=0$ and
$u_6=-180$ remain available in \eqref{long68:eq:low-channel-support}.
The least positive moment is then $115\cdot180/\gcd(180,55)=4140$.
Minimising the moment and minimising the largest supported index are
therefore different problems.

The factor beyond $L_D$ is not always $12$. At $D=6$, the same recurrence
gives
\[
 L_6=9839515,\qquad a_6=-2242555,\qquad G_6=840.
\]
To verify the last value, $u_7=0$, while $840=\lcm(1,\ldots,8)$ divides
every $u_n$ for $n\ge8$ by the preceding divisibility. Conversely,
$u_8=-4200$, $u_{12}=-14386680$ and $\gcd(u_8,u_{12})=840$.
Thus $\gcd(G_6,a_6)=35$ and the least positive moment is $\mu_6=24L_6$,
not $12L_6$. It is attained by $24K_6-16240U_8+U_{12}$: the coefficient
at index $1$ is zero and the coefficient at $12$ is $-1$. The
attainable-moment formula records restrictions that the uniform factor
$12$ alone does not capture.

We can now express the remainder using the basis coefficients in
\eqref{long68:eq:low-channel-basis}. This refines the integer-difference
identity at the start of \S\ref{long68:sec:channels} by giving its integer
term explicitly for the classified vectors.

\begin{theorem}[how the coefficient choices change the remainder]\label{long68:res:residual-transparency}
For the vector in \eqref{long68:eq:low-channel-basis},
\[
 \mathcal R\!\left(tK_D+\sum_{n>D}z_nU_n\right)
 =tL_D(\Ser-H_D)+\sum_{n>D}z_n.
\]
The residual series converges for every finite vector supported away from
index zero.  A zero-moment vector has integral residual, and any two finite
vectors with the same factorial moment have residuals differing by an integer.
\end{theorem}
\leannote{Lean: \lproof{ErdosProblems/Erdos68/PaperCompleteResidualIdentity.lean}{116}{residual_transparency}, \lproof{ErdosProblems/Erdos68/PaperCompleteResidualIdentity.lean}{166}{summable_fullResidual}, \lproof{ErdosProblems/Erdos68/PaperCompleteResidualIdentity.lean}{185}{zero_moment_residual_integral}, \lproof{ErdosProblems/Erdos68/PaperCompleteResidualIdentity.lean}{195}{equal_moment_residual_integer_difference}.}

\begin{proof}
For $d>D$, we have $V_d(K_D)=L_D$ and
$V_d(U_n)/(d!-1)=\mathbf1_{d=n}$. Summing these identities gives
$tL_D(S-H_D)+\sum_{n>D}z_n$, as claimed. The first series converges and
the second sum is finite.

Absolute convergence and the general integer-difference identity were
proved at the start of \S\ref{long68:sec:channels}, including for the
auxiliary index $1$. Setting the moment to zero proves integrality;
subtracting the identity for two vectors of equal moment proves the
last assertion.
\end{proof}

Formal counterparts are linked for
\href{https://github.com/wcook04/plectis-erdos/blob/f20f727ae44d6dc3b24036bf63118b8fe591912e/ErdosProblems/Erdos68/PaperCompleteMomentIdeal.lean#L242}{a primitive vector attaining the least positive moment}
(Theorem~\ref{long68:res:moment-ideal}) and
\href{https://github.com/wcook04/plectis-erdos/blob/f20f727ae44d6dc3b24036bf63118b8fe591912e/ErdosProblems/Erdos68/PaperCompleteResidualIdentity.lean#L195}{the integer difference between remainders with equal moments}
(Theorem~\ref{long68:res:residual-transparency}).

The factor $12$ used below follows directly on support $n\ge2$.
For $n=2,3$ one has $n!=2W_{2,n}$. For $n\ge4$, both $n!$ and
$2W_{2,n}$ are divisible by $12$: if $n=2k$ or $2k+1$, then
$W_{2,2k}=k!\prod_{j=1}^k(2j-1)$ is divisible by $6$ for $k\ge2$.
Thus $12\mid M(c)-2V_2(c)$. Every $d!-1$ is coprime to $6$, so
$V_2=\cdots=V_D=0$ implies $12L_D\mid M(c)$.
This factor is attained at $D=2$ and $D=3$ by
$-6e_2+e_4$ and $-6e_2-8e_3+5e_4$, respectively.
The depth-$6$ example above shows why it is not always the whole restriction.

The basis formula \eqref{long68:eq:low-channel-basis} describes all solutions
of the weighted-sum equations, and \eqref{long68:eq:low-channel-support}
imposes the support restriction. The remainder identity then determines
their residues modulo $\Z$. In particular, $\mathcal R(c)$ is an integer
linear form in $1$ and $S$, and its distance from the nearest integer equals
that of $M(c)S$. Integer translation can choose a representative near zero,
but cannot change this distance. To exclude every rational denominator this way, each positive integer $q$
must divide the moment of some nonintegral remainder in the family.
Eventual divisibility by every fixed $q$ is sufficient; growth of the
moments alone is not. For example, the compulsory factor $12L_D$ has
$3$-adic valuation exactly $1$ for every $D$, so this necessary divisor
alone never guarantees that $9$ divides the moment.

A truncation cutoff can exceed the largest supported index without changing
the vector. For example, $c=-6e_2+e_4$ has $M(c)=12$ and
$\mathcal R(c)=12S-17$. At cutoff $4$, its finite part is $-239/115$:
the gap $9/115$ is smaller than the tail bound $24/119$. At cutoff $5$,
the finite part is $-27061/13685$, whose gap $13376/13685$ exceeds the
new bound $24/719$. Thus the second cutoff certifies nonintegrality of
the same remainder, although the first does not. This excludes only
denominators dividing $12$, already covered by the finite exclusions.

For fixed $c$ with $M(c)>0$, the finite part increases to
$\mathcal R(c)$, with omitted tail $M(c)(S-H_N)$. The upper bound
$2M(c)/((N+1)!-1)$ is smaller than the finite part's gap to the next
integer at every sufficiently large cutoff if and only if
$\mathcal R(c)\notin\Z$. Indeed, for a nonintegral remainder the gap is
eventually at least
$\lfloor\mathcal R(c)\rfloor+1-\mathcal R(c)>0$, whereas the upper bound
tends to zero. For an integral remainder, the eventual gap equals the
omitted tail itself and is smaller than that upper bound. At any one
cutoff, success places the remainder strictly above the finite part and
below the least integer strictly exceeding it, so it implies nonintegrality.
The short note writes out the inequality in its section ``The remaining
real comparison''; increasing the cutoff alone does not prove that it succeeds. This is a truncation choice, distinct from prescribing
the largest nonzero coefficient index.

To impose this small-tail comparison at the largest supported index, one
needs more than $L_D\mid M$: the positive moment must also fit below a
factorial at that index. The next theorem combines those two demands at
$D=2t^2$. Its variable $R$ is a support parameter, not the remainder or
the modulus $R_p$ of the preceding section. If a larger truncation cutoff
is used instead, the same numerical bound applies to that cutoff, not
to the vector's actual support. The extra size assumption on $M$ is not
a consequence of cancelling the weighted sums.

\label{long68:res:lead-channel-radius}\phantomsection
\begin{theorem}[a lower bound for the support parameter]\label{long68:res:channel-radius}
Let $t,M,R\in\N$ satisfy
\[
 t\ge2^{32},\qquad M>0,\qquad \Lc_{2t^2}\mid M,\qquad M<(R+1)!-1 .
\]
Then $3t^3<2(R+1)$.  Consequently no family satisfying these hypotheses for
all sufficiently large $t$ has $(R(t)+1)/t^3\le3/2$ eventually, and none has
$R(t)=o(t^3)$.

\end{theorem}
\leannote{Lean: \lproof{ErdosProblems/Erdos68/PaperCompleteExisting.lean}{234}{square_subsequence_radius}, \lproof{ErdosProblems/Erdos68/PaperCompleteExisting.lean}{258}{radius_no_eventual_ratio_upper}, \lproof{ErdosProblems/Erdos68/PaperCompleteExisting.lean}{277}{radius_not_littleO}.}

\begin{proof}
Put $D=2t^2$, $k=2t$ and $u=D-k+1=2t^2-2t+1$.  Lemma~\ref{long68:res:segment} at
$N=D$, with $n!-1\ge n!/2$ and $\log n!\ge n\log n-n$, and with
$\Lc_D\le M<(R+1)!$, gives
\begin{equation}\label{long68:eq:radius-log}
 2t\bigl(u\log u-u-\log2\bigr)
 <(R+1)\log(R+1)+\binom{2t+1}{3}\log(2t^2).
\end{equation}


Suppose $R+1\le\tfrac32t^3$.  Since $t\ge16$ one has
$\tfrac{15}{8}t^2\le u\le2t^2$ and $\log u\ge2\log t$, so the left side of
\eqref{long68:eq:radius-log} is at least
$\tfrac{15}{2}t^3\log t-4t^3-2t\log2$.  Using
$\binom{2t+1}{3}<\tfrac43t^3$ and $\log\tfrac32<\log2$, the right side is at
most
\[
 \tfrac32t^3\bigl(\log2+3\log t\bigr)+\tfrac43t^3\bigl(\log2+2\log t\bigr).
\]
Dividing by $t^3$ and rearranging, \eqref{long68:eq:radius-log} would require
\[
 \tfrac13\log t<4+\frac{2\log2}{t^2}+\tfrac{17}6\log2 .
\]
But $\log t\ge32\log2$ and $\tfrac23<\log2<1$ give
$\tfrac13\log t-\tfrac{17}6\log2\ge\tfrac{47}6\log2>5$, while
$4+2\log2/t^2<5$.  This contradiction proves the finite inequality, and the
two sequence conclusions follow.
\end{proof}

These hypotheses are compatible: for a fixed $t$, one may take
$M=\Lc_{2t^2}$ and then choose $R$ large enough. The restriction is on how
small $R$ can be. Vanishing of $V_2,\ldots,V_{2t^2}$ forces the divisibility
condition, but not the upper bound $M<(R+1)!-1$. An application to a remainder
estimate must establish that additional size bound and $M>0$. The theorem
then rules out a parameter growing more slowly than cubically. It is a
bound on the actual support only when $R$ denotes its largest index and
the stated size condition holds there; it does not prove irrationality.

Theorem~\ref{long68:res:lcm-growth} raises the asymptotic constant in the same
estimate.

\begin{corollary}[the asymptotic lower bound]\label{long68:res:radius-constant}
Let $M(t),R(t)$ satisfy $M(t)>0$, $\Lc_{2t^2}\mid M(t)$ and
$M(t)<(R(t)+1)!-1$ for all sufficiently large $t$.  Then
\[
 \liminf_{t\to\infty}\frac{R(t)+1}{t^3}\ \ge\ \frac{16}{9}.
\]
\end{corollary}
\leannote{Lean: \lproof{ErdosProblems/Erdos68/PaperCompleteLiminf.lean}{54}{asymptotic_radius_constant_liminf}.}

\begin{proof}
Put $r=R(t)+1$, so $\Lc_{2t^2}\le M(t)<r!$ and $\log \Lc_{2t^2}<r\log r$.
Since
\[
\frac{(2t^2)^{3/2}\log(2t^2)}{t^3\log t}\longrightarrow4\sqrt2,
\]
Theorem~\ref{long68:res:lcm-growth} gives, for every $\varepsilon>0$ and
all sufficiently large $t$,
\[
 \log \Lc_{2t^2}\ge
 \Bigl(\tfrac{16}{3}-\varepsilon\Bigr)t^3\log t.
\]
Suppose $r\le ct^3$ for arbitrarily large $t$, with $0<c<16/9$ fixed.  On that
subsequence $r\log r\le ct^3(3\log t+\log c)=\bigl(3c+o(1)\bigr)t^3\log t$.
Choosing $0<\varepsilon<16/3-3c$ contradicts the lower bound for
sufficiently large $t$.
\end{proof}

This corollary does not assert the strict inequality
$R+1>\tfrac{16}9t^3$ eventually. Indeed, the finite logarithmic constraint
\eqref{long68:eq:radius-log} allows equality at that constant. To check
this directly, set $R+1=16t^3/9$ with $t\ge4$. Since
$u=2t^2-2t+1\le2t^2$ and $x\log x-x$ is increasing for $x\ge1$, the
left side of \eqref{long68:eq:radius-log} is at most
$4t^3(\log(2t^2)-1)-2t\log2$. Subtracting this bound from the right side gives
\[
 t^3\left(4+\frac{16}{9}\log\frac{16}{9}-\frac83\log2\right)
 +t\left(\frac53\log2-\frac23\log t\right)
 >\frac43t^3-\frac23t^2>0,
\]
where $\log2<1$ and $\log t\le t$ suffice for the last comparison.
Integer pairs with $9(R+1)=16t^3$ exist for every multiple $t\ge6$ of $3$.
This verifies the numerical constraint, not the existence of a coefficient
vector with that support. The asymptotic assertion in
Corollary~\ref{long68:res:radius-constant} instead follows from
Theorem~\ref{long68:res:lcm-growth}; the finite-block Lean source does not
carry that deduction. The corresponding Lean source statement is
\href{https://github.com/wcook04/plectis-erdos/blob/0b500c7cf8e8bb7ae343484378df02f277fb8194/lean/ErdosProblems/Erdos68/PaperCompleteLiminf.lean#L54}{the asymptotic lower bound for the support parameter}.

At a prime index, the two-term vector $T_p$ already changes only one
weighted sum; there are no proper divisors to eliminate.

\begin{theorem}[changing just one weighted sum]\label{long68:res:translator}
Let $p\ge3$ be prime and let $c_{p-1}=p$, $c_p=-1$, with every other coefficient
zero.  Then $M(c)=0$, $\chn pc=p!-1$, and $\chn dc=0$ for every $d\ge2$ with
$d\ne p$.
\end{theorem}
\leannote{Lean: \lproof{ErdosProblems/Erdos68/PaperCompleteExisting.lean}{299}{prime_channel_corrector}.}

\begin{proof}
Since $p\ge3$, both indices $p-1$ and $p$ lie in the prescribed coefficient
domain $i\ge2$.  The moment is $p\,(p-1)!-p!=0$.  For $2\le d<p$ the quotient identity
$\lfloor(p-1)/d\rfloor=\lfloor p/d\rfloor$ holds, since $d\nmid p$, so the
two factorial weights carry the same power of $d!$ and the weighted sum evaluates to
$p\,(p-1)!/(d!)^{\lfloor p/d\rfloor}-p!/(d!)^{\lfloor p/d\rfloor}=0$.  For
$d>p$ both indices lie below $d$, so both weights are the plain factorials
and the same cancellation occurs.  At $d=p$ the weights are $(p-1)!$ and
$p!/p!=1$, giving $p\,(p-1)!-1=p!-1$.
\end{proof}

For $p=2$, the same auxiliary identities would require $c_1=2$.  They do
not define an admissible vector on the domain $i\ge2$, since every vector
supported on that domain has $c_1=0$.

\noindent Adding an integer multiple of this vector changes $V_p$ by
that multiple of $p!-1$, without changing $M$ or any other $V_d$.
Suppose $c$ has nonzero moment, $V_2(c)=\cdots=V_D(c)=0$, and every
supported index exceeds a prescribed integer $B$. Choose a prime
$p>\max(D,B+1,2)$ and replace $c$ by
\[
 c-\lfloor\mathcal R(c)+1/2\rfloor T_p.
\]
Both new indices, $p-1$ and $p$, exceed $B$. Since $p>D$, the prescribed
weighted sums remain zero, and the moment is unchanged. The new remainder is
$\mathcal R(c)-\lfloor\mathcal R(c)+1/2\rfloor\in[-1/2,1/2)$.
The progression construction below supplies such an initial vector:
choose its first index $r>B$.

This rounding proves a size bound, not nonvanishing: an integral remainder
becomes zero. In fact, a rounded remainder is nonzero exactly when the
original remainder is nonintegral. Choosing larger adjustment primes
cannot extend the denominator coverage, since the moment stays fixed.
The chosen integer
$\lfloor\mathcal R(c)+1/2\rfloor$ depends on the remainder itself. The
formula establishes the existence of a representative in the stated
interval; it does not independently certify that this representative is
nonzero.

\paragraph{A primitive solution on an arithmetic progression.}
To extend the example $-6e_2+e_4$ to cancellation at larger $D$, place
the support on an arithmetic progression with its first index prescribed.
Let $D,r\ge2$ and let $\ell$ be a positive multiple of $2,\ldots,D$. Equal
spacing makes $V_d=0$ a root condition at $(d!)^{-\ell/d}$. We choose
the degree-$D-1$ polynomial with exactly these roots and scale its
coefficients so that the last supported coefficient is $1$. Put
\[
i_j=r+j\ell\ (0\le j<D),\quad N=r+(D-1)\ell,\quad
\alpha_d=(d!)^{\ell/d},
\]
\[
\prod_{d=2}^D(\alpha_dX-1)=\sum_{j=0}^{D-1}h_jX^j,
\qquad A=\prod_{d=2}^D\alpha_d.
\]
At index $i_j$ put the coefficient
\[
c_{i_j}=\frac{N!h_j}{A i_j!},
\]
and put every other coefficient equal to zero. We verify that these
coefficients are integers and that $V_2=\cdots=V_D=0$.
For $2\le d\le D$, divisibility of $\ell$ by $d$ gives
$\lfloor i_j/d\rfloor=\lfloor r/d\rfloor+j\ell/d$. Hence
\[
V_d(c)=
\frac{N!}{A(d!)^{\lfloor r/d\rfloor}}
\sum_{j=0}^{D-1}h_j\alpha_d^{-j}=0,
\]
because $\alpha_d^{-1}$ is a root of the displayed polynomial. Evaluating
that polynomial at $1$ gives the moment
\[
M=N!\prod_{d=2}^D(1-1/\alpha_d)>0.
\]

For integrality, each term of $h_j/A$ is, up to sign, the reciprocal of
a product of $D-1-j$ distinct factors $\alpha_d$. In that denominator,
each selected $d$ contributes $\ell/d$ blocks of size $d$. The total size is
$(D-1-j)\ell=N-i_j$, so the multinomial coefficient shows that this
denominator divides $N!/i_j!$. Thus every $c_{i_j}$ is an integer.
Since $h_{D-1}=A$, the coefficient at $N$ is $1$, and the coefficients
have gcd one. This does not establish minimal support. The smallest
permitted step is $\ell=\operatorname{lcm}(2,\ldots,D)$; taking $D=r=\ell=2$
recovers $-6e_2+e_4$.

The stronger divisibility
\[
r!\prod_{d=2}^D(\ell/d)!\prod_{d=2}^D(\alpha_d-1)\mid M
\]
follows because
\[
\frac{N!}{r!\,A\prod_{d=2}^D(\ell/d)!}
\]
is an integer: it counts partitions into one distinguished block of size
$r$ and, for each $d$, an unordered collection of $\ell/d$ blocks of size
$d$. Their sizes sum to $N$. Using
$M=(N!/A)\prod_{d=2}^D(\alpha_d-1)$ now gives the asserted divisibility.
Put $v=\max(r,\ell/2)$, an integer since $2\mid\ell$. The displayed
factor divides $M$, so $v!\mid M$. Since $D\le\ell\le2v$ and $r\le v$,
\[
 N=r+(D-1)\ell\le4v^2-v<4v^2,
 \qquad v\ge\lfloor\sqrt N/2\rfloor+1.
\]
Consequently $(\lfloor\sqrt N/2\rfloor+1)!\mid M$. Every positive integer
$q\le\lfloor\sqrt N/2\rfloor+1$ therefore divides $M$, as does every
divisor of that factorial. In particular, $N\ge4(q-1)^2$ suffices for
$q\mid M$. This includes denominators growing with $N$.

The construction also supplies the required absolute smallness of the
omitted tail at its own support endpoint. Every factor $1-1/\alpha_d$
lies strictly between $0$ and $1$, so $0<M<N!$ and
\[
 0<M(S-H_N)<\frac{2M}{(N+1)!-1}<\frac2N,
\]
where the last inequality follows from $(N+1)!-1>NN!$ for $N\ge2$.
In particular, the moment obeys $M<(N+1)!-1$ without an additional
hypothesis for this family. What is not established is that the displayed
upper bound is smaller than the finite part's gap to the next integer.
That gap depends on the same vector and may shrink with $N$; an upper
bound tending to zero does not supply the required comparison. Neither
this size estimate nor the divisibility proves nonintegrality of the
remainders.

\section{Finite denominator exclusions}\label{long68:sec:finite}
\label{long68:res:lead-denominator-exclusions}\phantomsection

Two finite computations constrain a hypothetical denominator $q$ in
$\Ser=a/q$, and they constrain different features of it:
\begin{equation}\label{long68:eq:finite-bounds}
 q\nmid299999!,\qquad q\ge2^{39990}>10^{12038}.
\end{equation}
The size bound holds for every integer $a$ and positive natural $q$ with
$\Ser=a/q$; no reducedness assumption is needed. Its evidence is the exact
continued-fraction computation described below. The stronger candidate
\href{https://github.com/wcook04/plectis-erdos/blob/27c2fc5fe3c55fe547feaeb4c9bb68b3cf63a5bf/lean/ErdosProblems/Erdos68/PaperCompleteFiniteSizeCertificate.lean#L58}{finite-size Lean certificate}
is not cited as kernel verification: its source records that the computation
has not yet been kernel checked. The factorial-grid exclusion
$q\nmid299999!$ is supported separately by the external GMP computation;
its full Lean certificate also remains an outstanding formal obligation.

Neither condition implies the other. The prime $300007$
fails to divide $299999!$ but is smaller than $2^{39990}$, whereas $299999!$
satisfies the size lower bound but fails the nondivisibility condition.
High powers of small primes further show why nondivisibility is not a
smoothness exclusion.

The first exclusion comes from an exact interval carry census.  The supplied record describes an
independently implemented GMP integer computation certifying all $299998$
carry cells for $3\le m\le300000$ with a scale of $2^{5025679}$ and $96$
guard bits, using no floating-point arithmetic; at every step the outward
interval is proved to lie inside one half-open unit cell before the carry is
recorded.  Its unit carries occur exactly at
\[
 52,\ 591,\ 1030,\ 1407,\ 1438,\ 2164,\ 4258,\ 10991,\ 21236,
\]
so $b_{300000}\ne1$.  Theorem~\ref{long68:res:carry-equivalence} then gives
$q\nmid299999!$ and in particular $q\ge300000$. The full-range carry replay reported in the supplied manuscript was not
repeated for this prose revision. A separately recorded replay through
$m=4000$ reproduces the unit-carry prefix $52,591,1030,1407,1438,2164$ with
a non-unit carry at the endpoint, giving $q\ge4000$ on its own. These are
external computations; the implication from a non-unit carry to a
factorial-divisibility exclusion is the kernel-checked statement.

The first exclusion says that the least positive integer $r$ with
$q\mid r!$ is at least $300000$. This is the factorial index used by Sondow
\cite[\S3, Theorem~1]{sondow2006}, not the largest prime factor of $q$.
No approximation bound for $e$ is being applied to $\Ser$.

The second exclusion has a short integer description.  Put $D=2^{80000}$ and
let $N$ be the first integer with $N!-1>D$.  Define
\[
 \ell=\sum_{n=2}^{N-1}\left\lfloor\frac D{n!-1}\right\rfloor,
 \qquad
 u=\ell+(N-2)+\left\lfloor\frac{2D}{N!-1}\right\rfloor+1 .
\]
Each rounded prefix term loses less than one. Since
$(n+1)!-1>(n+1)(n!-1)$, successive terms in the tail from $N$ have ratio
less than $1/(N+1)$. Their sum is consequently smaller than
$(N+1)/(N(N!-1))<2/(N!-1)$, proving $\ell/D<\Ser<u/D$.

There are exactly $N-2$ rounded prefix terms.  The extra $+1$ in $u$
keeps the upper enclosure strict even if $2D/(N!-1)$ happens to be an
integer; the positive omitted tail keeps the lower enclosure strict.  Here
$N=7054$ and $u-\ell=7053$.  Apply the continued-fraction algorithm to both
endpoints at once, retaining a quotient only while the integer parts agree
and both remainders are nonzero, and reversing the endpoint order at each
inversion.  There are $23449$ common quotients, and the convergent
denominators
\[
 Q_{-2}=1,\quad Q_{-1}=0,\quad Q_j=a_jQ_{j-1}+Q_{j-2}
\]
satisfy $Q_{23448}\ge2^{39990}$.  Every rational in the open enclosure shares
this initial segment, so its reduced denominator is at least $Q_{23448}$.
Indeed, write $P_j$ for the convergent numerators.  Such a rational either
equals $P_{23448}/Q_{23448}$ or has a complete quotient $A/B>1$ after the
segment, with $A$ and $B$ coprime, and then equals
$(P_{23448}A+P_{23447}B)/(Q_{23448}A+Q_{23447}B)$; the determinant identity
$P_jQ_{j-1}-P_{j-1}Q_j=\pm1$ makes both fractions reduced
\cite[\S1.12(ii), (1.12.5)--(1.12.7), (1.12.20)--(1.12.21)]{nist-dlmf}.  The
argument is about a common prefix and does not require the last convergent to
lie inside the enclosure.  The integer comparison $2^{39990}>10^{12038}$
closes \eqref{long68:eq:finite-bounds}.  The
\href{https://github.com/wcook04/plectis-erdos/blob/cd3bd6fe245867a435e15d503ccedd699c5d02e2/scripts/check_erdos68_continued_fraction.py}{integer
replay} and its
\href{https://github.com/wcook04/plectis-erdos/blob/cd3bd6fe245867a435e15d503ccedd699c5d02e2/verification/erdos68-continued-fraction.json}{receipt}
record scale $80000$ bits, bracket width $7053$, $23449$ certified partial
quotients, power-of-two exponent $39990$ and strict decimal exponent $12038$.


\paragraph{Carry computation record.}
The receipt \path{verification/erdos68-strict-successor.json}, with schema
\path{erdos68_strict_successor_exact_interval_v1}, records the scale and guard
bits above, the driver and backend used for its run, the event-trace digest,
and the final enclosure. The source packet names the driver
\path{scripts/check_erdos68_strict_successor.py} and the backend
\path{scripts/check_erdos68_strict_successor_gmp.cpp}. Their published digests,
followed by the receipt payload digest, are
\begin{quote}\small
Driver: \path{f25b5bd8ffbfd66bdbb42ca6a07a936aab9b4eb7d3a57446725be55acb642191}.\\
Backend: \path{43308ddc3902dd537b789ae6e0648054a0d752adf07e5d6e1d829811e9342083}.\\
Payload: \path{044ef27a60f7294c39b8279b7f04e22e1234a942e2862e9e8c64eb66505af29b}.
\end{quote}
The supplied record states that the event trace and unit-carry certificate
match the previously retained receipt. These identifiers distinguish the
reported full run from the separate computation through $4000$; they do
not certify a new replay in this revision.

Both certificates are finite. A further computation can enlarge an
exclusion, but need not do so: a narrower enclosure may retain the same
common continued-fraction prefix, and a later unit carry supplies no new
carry exclusion. Neither finite calculation excludes an eventual
unit-carry tail.  The continued-fraction identities are classical
\cite[\S1.12(ii)]{nist-dlmf}; the enclosure, the prefix length and the
exponent are outputs of the computation above.

\section{The remaining arithmetic inputs}\label{long68:sec:open}

The exact target is $m\nmid Z_m$ at arbitrarily large indices, as in
\eqref{long68:eq:strict-misses}; the equivalent comparison with fractional
parts is \eqref{long68:eq:lower-escape}. The estimates discussed below
would suffice for this target but are not proved here. The finite test
\eqref{long68:eq:finite-escape} implies escape at each index; its occurrence
at arbitrarily large integer indices is also equivalent to irrationality,
as proved in \S\ref{long68:sec:carry}. Thus the missing step is to prove
that this actual rational sequence meets the test at arbitrarily large
indices, not to supply another equivalence. Restricting to a prime subsequence or replacing
a numerator by a residue imposes additional conditions. A model obeying
only selected congruences is not the actual sequence.

\paragraph{Value sets do not determine the weighted residue sum.}
Garaev, Luca and Shparlinski's harmonic-sum estimate
\cite[Theorem~1 and (5)]{garaev-luca-shparlinski-harmonic} concerns
unconditioned harmonic sums, not sums restricted to the indices where
$n!\equiv1\pmod q$.
The value-set literature addresses different invariants.  Banks, Luca,
Shparlinski and Stichtenoth \cite{banks-et-al2005} provide historical
missing-value context; Klurman and Munsch
\cite[Theorems~2.1--2.2]{klurman-munsch2017} distinguish unconditional
averages from GRH-dependent improvements.  Grebennikov, Sagdeev,
Semchankau and Vasilevskii
\cite[Corollaries~1.2 and~1.4]{grebennikov-et-al2024} study sizes of value
and quotient sets, as do the factorial-residue and representation results
of Hu \cite{hu2026} and Garaev and Pardo \cite{garaev-pardo2026}.
None is a noncancellation theorem for the reciprocal unit weights here.
A lower bound for the number of distinct residues does give an upper
bound on the multiplicity of residue $1$, since every other value needs
at least one index. But that bound does not determine the reciprocal
unit weights or test whether they cancel modulo
$q$. The two explicit cancellations in \S\ref{long68:sec:prime-powers} illustrate
why the weights matter.
Hu's September preprint \cite[Theorem~3.1]{hu-september2026} considers
pairwise distinct maps $T_\theta(x)=b_\theta+c_\theta/x$ over $\mathbb F_p$,
with $c_\theta\ne0$. An internal transition is a pair $(x,\theta)$ for which
$x\in A\smallsetminus\{0\}$ and $T_\theta(x)\in A$.
With at least $M$ distinct such pairs, at most $|A|$ parameters, and a
bound $\mu\ge1$ on the number of ordered parameter pairs representing each
nonidentity map $T_\eta T_\theta^{-1}$, the theorem gives $|A|\gg\min\{M,p\}^{8/15}\mu^{-4/15}$. Its incidence estimate comes
from Stevens and de Zeeuw \cite[Theorem~4]{stevens-dezeeuw2017}.
To apply this to the residue classes of indices satisfying
$n!\equiv1\pmod q$, one would first have to specify the set $A$ and the
transition maps. The number of such indices is not automatically a
lower bound for $M$: different indices may produce the same pair
$(x,\theta)$, and their factorial values are all $1$. One must prove that
enough distinct transitions remain, as well as bound the ordered-pair
multiplicity $\mu$. The unconditioned value-set theorem does not provide
these additional facts.

\paragraph{The shared denominator and the gap must be controlled together.}
For the blocks selected from first prime occurrences, the second inequality
in \eqref{long68:eq:factor-split} is
\begin{equation}\label{long68:eq:weighted-target}
 \log\widetilde C_p
 <\log\left(\frac{2p^2}{2p+1}\cdot
       \frac{\prod_{j=p}^{2p-1}j}{q}\right).
\end{equation}
The first inequality is \eqref{long68:eq:global-scale}. Both must hold at
the same arbitrarily large parameters.

The order of quantifiers is part of the proposed result: for every bound,
one parameter must satisfy both inequalities. Two unrelated infinite
subsequences would not prove the combined criterion.

For a prime $r$ and an integer $e\ge1$, let
$h_{r,e}(p)=\#\{i\in I_p:r^e\mid i!-1\}$. The exponent of $r$ in the
shared part is at least $e$ exactly when $h_{r,e}(p)>1$.
Removing the exponent already present in $F_p$ therefore gives
\[
 v_r(\widetilde C_p)
 =\#\{e\ge1:h_{r,e+v_r(F_p)}(p)>1\}.
\]
When $r\nmid F_p$, this counts the positive exponents for which two or
more denominators are divisible by $r^e$. The spacing bound
$h_{r,e}(p)(e+1)\le2p+e-2$ follows from the gap-product argument:
if $i<j$ are two such indices, then $j<r$ and
$0<r^e\le j!/i!-1<r^{j-i}$, so $j-i\ge e+1$. All indices lie between
$2$ and $2p-1$, giving the bound; the zero-hit case is immediate.
To estimate $\log\widetilde C_p$, these exponent counts must be multiplied
by $\log r$ and summed over primes. A count ignoring the sizes of the
primes does not give \eqref{long68:eq:weighted-target}. For a prime $q$,
Wilson reflection of an odd index $n<q$ contributes a repeated divisor
only when the reflected index is distinct and in the same block, with
both indices at least $2$. The endpoint $n=q-2$ reflects to $1$ and
contributes no such pair, as shown in \S\ref{long68:sec:prime-powers}.

Even for two admissible indices, the reflection identity is only modulo
$q$ \cite[(4), p.~462]{stewart2004}. For example, exact multiplication gives
\[
 609!\equiv1\pmod{971^2},\qquad
 361!\equiv736019=1+758\cdot971\pmod{971^2},
\]
where $361=971-609-1$. Both indices lie in $I_{306}$, but their factorial
denominators have $971$-exponents at least $2$ and exactly $1$, respectively.
Thus reflection supplies two hits modulo $q$, not necessarily modulo
$q^e$ for $e>1$. Each level in the valuation sum requires its own
congruence; neither distinctness nor a higher exponent may be inferred
from the prime-level identity alone.
Theorem~12 of Garaev, Luca and Shparlinski
\cite[arXiv v1, Thm.~12, p.~16]{garaev-luca-shparlinski} is a multiplicity
bound and supplies no lower bound for the distance to the next integer in
\eqref{long68:eq:global-scale}.

For the first inequality alone, a joint mean bound
\[
 \frac1{|I_X|}\sum_{p\in I_X}
 \left(\log\widetilde C_p-\log\frac{\rho_p}{R_p}\right)
 <(1-\varepsilon)X\log X
\]
on nonempty sets $I_X\subseteq[X,2X]\cap\mathbb N$, with $R_p>1$, a fixed
$0<\varepsilon<1$, and arbitrarily large $X$, would suffice: at least one
summand is at most the average, while the required upper bound is
$\log(K_p/((2p+1)F_p))=p\log p+O(p)$, uniformly for $p\in[X,2X]$.
Two separate estimates on unrelated sets of indices would not give this
joint comparison.

\paragraph{What two gap-product tests do not show.}
For integers $h\ge1$ and $x\ge0$, let
$F_h(X)=\prod_{j=1}^h(X+j)$. Suppose $x+h<q$ for a prime $q$.
If $x!$ and $(x+h)!$ both equal $1$ modulo $q$ and the harmonic sums
through these indices are equal modulo $q$,
then $F_h(x)=1$ and
$F'_h(x)=F_h(x)\sum_{j=1}^h(x+j)^{-1}=0$ in that field.  An exact calculation gives
\[
 \gcd(F_{12}-1,F'_{12})
 =(X-4626)(X-7848)\quad\hbox{in }\mathbb F_{12487}[X].
\]
Thus a uniform bound of one repeated root is false. But $4626!\equiv442$ and $7848!\equiv6300$
modulo $12487$, so neither root also satisfies $x!\equiv1$. Counting
repeated roots of $F_h-1$ therefore does not count the indices required
by the factorial problem. These are finite modular calculations, not a
bound on the number of indices satisfying both congruences.

Stewart's gap-product estimate
\cite[Lemma~2, (21)--(24)]{stewart2004} requires the two products to be
unequal. He supplies that step separately: in the proof of his bound
(14), a prime-distribution argument separates the products when the common
divisor is large; the smaller-divisor case is immediate
\cite[p.~467, proof of Lemma~2]{stewart2004}.

The non-power theorem gives a different argument when the gaps are
adjacent. For integers $2\le i<j<k$, equality $j!/i!=k!/j!$ would make $k!/i!$ a
square, contrary to that theorem \cite[Theorem~1]{erdos-selfridge1975}.
Selecting two short gaps with nearly equal lengths need not preserve
adjacency, so this observation alone cannot replace Stewart's nonvanishing
argument. The wider equal-product problem has a separate history
\cite{macleod-barrodale1970}; the doubled-length theorem discussed by
Nguyen Xuan Tho \cite[Theorem~1]{nguyen-tho2026} is another special case,
not arbitrary near-equal-length nonvanishing. No stronger lcm exponent is
inferred from these local observations.

\paragraph{Prime powers already seen at earlier indices.}
Recall that $\Delta_n=P^\Delta_n/Q^\Delta_n$ is the positive gap in lowest terms. Define
\[
 \mathcal A_n=\prod_{\substack{r\ \mathrm{prime},\ r\mid n!-1\\
     r\mid k!-1\ \text{for some }2\le k<n\\
     v_r(Q^\Delta_n)<v_r(n!-1)}}r^{\,v_r(n!-1)} .
\]
Thus the product includes the full power of a prime in $n!-1$ when that
prime has appeared before but has smaller exponent in $Q^\Delta_n$.
The recurrence for the gaps is
\[
 \Delta_{n+1}=n\Delta_n-b_n-\frac1{n!-1}.
\]
For a prime $r$ in the product, put $e=v_r(n!-1)$. The reduced denominator
of $n\Delta_n-b_n$ has $r$-exponent at most $v_r(Q^\Delta_n)<e$, whereas
$1/(n!-1)$ has exponent $e$. Clearing denominators leaves exactly one
numerator term nonzero modulo $r$, so the reduced denominator of
$\Delta_{n+1}$ has exponent exactly $e$. Hence
$\mathcal A_n\mid Q^\Delta_{n+1}$ and, when $\mathcal A_n>1$,
$P^\Delta_{n+1}\not\equiv0\pmod{\mathcal A_n}$.
The earlier occurrence of $r$ is not used in this divisibility argument;
it restricts the product to the primes proposed for this approach.
A sufficient quantitative estimate, at arbitrarily large indices $n$,
with $m=n+1$ and $d=\mathcal A_n>1$, is
\begin{equation}\label{long68:eq:amplification-target}
 \bigl((m+2)m!-2\bigr)Q^\Delta_m
 \le m^2(m!-1)\,(P^\Delta_m\bmod d).
\end{equation}
It implies the upper alternative of \eqref{long68:eq:finite-escape}. A nonzero
least representative can equal one even when the modulus is large.
Neither a lower bound for $\log\mathcal A_n$ nor infinitely many increases
in a previously occurring prime's exponent therefore proves
\eqref{long68:eq:amplification-target}: the least representative itself
must satisfy the displayed inequality. Repeated prime squares
can first occur before the Wilson index $r-2$: for example, $971$ divides
$361!-1$ and its square first divides $609!-1$, with $609<971-2$. Such an
example rules out confinement to that last index but supplies no unbounded
family satisfying the required inequality.

\paragraph{The upper alternative in the finite test.}
It would suffice for \eqref{long68:eq:finite-escape} to hold at arbitrarily
large indices; an unbounded set of prime indices would suffice as well.  The upper alternative is exactly
$\bigl((m+2)m!-2\bigr)Q^\Delta_m\le m^2(m!-1)P^\Delta_m$.  Replacing $P^\Delta_m$ by its least
nonnegative residue modulo a specified divisor of $Q^\Delta_m$ gives a stronger
sufficient test. No bound forcing this inequality at arbitrarily large
indices is established here. The full recurrence for the actual scaled
partial sums, including its initial value, is fixed; a rational example
for a weakened recurrence does not satisfy that full specification
(\S\ref{long68:sec:ext-nogo}). Nor does a zero canonical digit assert
membership in the specified real interval.

\paragraph{Conditions at twice a prime.}
For an odd prime $p$, the carry range and $Z_{2p}=2pZ_{2p-1}+1-b_{2p}$ give
\begin{equation}\label{long68:eq:doubled-prime}
 p^2\mid Z_{2p}
 \quad\Longleftrightarrow\quad
 \begin{cases}
 b_{2p}=1\ \hbox{and}\ p\mid Z_{2p-1},\ \hbox{or}\\
 b_{2p}=1+p\ \hbox{and}\ p\mid2Z_{2p-1}-1 .
 \end{cases}
\end{equation}


Reduction modulo $p$ forces $b_{2p}\equiv1\pmod p$, and the two displayed
values are the only possibilities in $[-1,2p-1]$; division by $p$ gives the
predecessor conditions.  For rational $\Ser=a/q$ and $p>q$, the tail estimate gives
$Z_{2p}=(2p)!\,\Ser$, which is divisible by $p^2$. A proof through
$p^2\nmid Z_{2p}$ must therefore exclude both displayed alternatives at
the same prime, for arbitrarily large primes.
That is a requirement of this particular test, not of every irrationality
argument. Rationality also forces $b_{2p}=1$ and
$Z_{2p-1}=(2p-1)!a/q$, with $p\mid Z_{2p-1}$: here $q\mid(p-1)!$,
and the remaining factorial factors include $p$. Thus either
$b_{2p}\ne1$ at arbitrarily large odd primes or
$p\nmid Z_{2p-1}$ at arbitrarily large odd primes would already contradict
rationality. The first is a special case of
Theorem~\ref{long68:res:carry-equivalence}; the second uses the same
factorial-clearing argument. Neither unbounded occurrence is proved here.

\paragraph{The cofactor form of the progression construction.}
For a positive starting parameter, a Vandermonde matrix gives another
formula for the progression vectors of \S\ref{long68:sec:channels}; it
does not resolve the remaining nonintegrality problem. For integers $n,t\ge0$, put
$s_n=((n+2)!)^2$ and $i_j=(t+j)s_n$ for $0\le j\le n+1$.
This step size matches the recorded construction; the determinant argument
only needs a positive step divisible by every $d=2,\ldots,n+2$.
Here the factorial-weight formula is also allowed at index $0$; the case
$t=0$ will show why a lower support condition is necessary. Form the
integer matrix $A$ with moment row $(i_j!)_j$ and weighted-sum rows
$(W_{d,i_j})_j$ for $2\le d\le n+2$.

To see that $\det A\ne0$, divide column $j$ by $i_j!$.
Since $d\mid s_n$, the $d$th weighted-sum row becomes
$((d!)^{-(t+j)s_n/d})_j$. Removing the nonzero factor
$(d!)^{-t s_n/d}$ from that row leaves a Vandermonde matrix with nodes
\[
1,\ (2!)^{-s_n/2},\ \ldots,\ ((n+2)!)^{-s_n/(n+2)}.
\]
These nodes are distinct: $(d!)^{1/d}$ is strictly increasing, since
$d!<(d+1)^d$. The Vandermonde determinant is the product of their
pairwise differences \cite[(1.3.13)]{nist-vandermonde}, and is therefore nonzero. None of the row or
column factors removed is zero, so the original determinant is nonzero
too.

Let $c$ be the cofactor vector of the moment row, and let $N_d$ be the
determinant obtained by replacing that row with $(W_{d,i_j})_j$.
Cofactor expansion gives $M(c)=\det A\ne0$ and $\chn dc=N_d$.
For $2\le d\le n+2$ the replacement duplicates a row, so $\chn dc=0$.  By
\eqref{long68:eq:channel-congruence} each $(\chn dc-M(c))/(d!-1)$ is an integer and
vanishes for $d>\max_ji_j$, so
\[
 \mathcal R_{n,t}:=\sum_{d>n+2}\frac{N_d}{d!-1},\qquad
 \mathcal R_{n,t}-\det(A)\,\Ser\in\Z .
\]
A family with $\min_ji_j=t\,s_n\to\infty$ and
$\mathcal R_{n,t}\notin\Z$ would prove irrationality. Indeed, for a fixed
denominator $q$, every entry of the moment row is divisible by $q$ once
the least support index is at least $q$, so $q\mid\det A$. If $S=a/q$,
the integer-difference identity would then force $\mathcal R_{n,t}\in\Z$.
Conversely, if $S$ is irrational, that same identity and $\det A\ne0$
make every $\mathcal R_{n,t}$ irrational. Taking, for example, $n=0$
and $t\to\infty$ gives the required family. Thus nonintegrality along
such a family is equivalent to the original irrationality question,
not a weaker conjecture obtained from the construction.

Merely making $(n,t)$ unbounded does not ensure that the support tends to
infinity: when $t=0$, the least support index is always zero. The coefficient
at $i_0$ really is nonzero: its cofactor uses columns $1,\ldots,n+1$, and
after nonzero row and column scalings it is a Vandermonde determinant
on the weighted-sum nodes. The eventual
identity $N_d=\det A$ determines the tail, but not the finite intermediate
sum, whose terms may have either sign. The construction solves the linear equations. What is still needed here
is a proof that the particular sum is nonintegral, whether by a gap
estimate or another arithmetic argument.

For $t\ge1$, primitive normalization identifies this vector exactly.
Take $D=n+2$, $\ell=s_n$ and $r=t s_n$ in the progression construction
of \S\ref{long68:sec:channels}, and put $N=r+(D-1)\ell$.
For any coefficient vector supported on these $D$ indices, the equations
$V_2=\cdots=V_D=0$ say that the polynomial
\[
 \sum_{j=0}^{D-1}c_{i_j}i_j!X^j
\]
vanishes at the $D-1$ distinct points
$\alpha_d^{-1}=(d!)^{-\ell/d}$, $2\le d\le D$.
Its degree is at most $D-1$, so it is a scalar multiple of
$\prod_{d=2}^D(\alpha_dX-1)$. Thus the rational solution space is
one-dimensional. The progression vector has coefficient $1$ at $N$,
so every integer solution is its integer multiple, with multiplier $c_N$.
In particular, the cofactor gcd is $|c_N|$, and the primitive cofactor
vector is the progression vector up to sign. With the sign chosen so
that its moment is positive, that moment equals
\[
 N!\prod_{d=2}^D\bigl(1-(d!)^{-\ell/d}\bigr).
\]
For example, $n=0,t=1$ gives support $\{4,8\}$, cofactor vector
$2520e_4-6e_8$, and positive-moment primitive vector $-420e_4+e_8$.

The factorial divisibility proved for the progression moment therefore
also applies to this primitive vector. Normalization still leaves the
nonintegrality condition equivalent to irrationality along families
whose least support index tends to infinity: the moment is nonzero,
every factorial below that index divides it, and
\eqref{long68:eq:channel-congruence} gives the integer-difference identity.
This identification does not estimate the distance from an integer.
The case $t=0$ is excluded from the identification with the stated
progression construction, whose indices are at least $2$.

\paragraph{Criteria from the literature that do not apply.}
Duverney's Theorem~3.1 assumes, among other conditions, quadratic growth
$cu_n^2\le u_{n+1}\le c'u_n^2$ for positive constants $c,c'$. For
$u_n=n!-1$, however, $u_{n+1}/u_n^2\to0$. His Corollary~3.2 assumes
convergence of the signed series $\sum_n(u_{n+1}/u_n^2-1)$ in (3.6).
Here its terms tend to $-1$, so that hypothesis fails as well
\cite[pp.~275, 285--287]{duverney}.  For a single denominator sequence
the rapid-growth criterion goes back to Erd\H{o}s
\cite[Theorem~1, p.~1]{erdos1975}; Barreto, Kang, Kim, Kova\v{c} and Zhang
treat products of consecutive denominators and weighted extensions
\cite[Thms.~2--3 and Rem.~4]{barreto-et-al}.  For $a_n=n!-1$,
$\log(n!-1)=O(n\log n)$ makes $(n!-1)^{1/\psi^n}\to1$ for every fixed
$\psi>1$, so neither Erd\H{o}s's hypothesis
$\limsup_na_n^{1/2^n}=\infty$ nor the growth hypotheses of their Theorems~2
and~3 hold.  These criteria do not decide $\Ser$; the useful content is the
target these papers identify: a subsequence with clearing integers small enough for the corresponding
remainders, or enough exact cancellation in those integers.
Theorem~\ref{long68:res:lcm-growth} measures the cost of clearing all
summands before reduction; it does not bound cancellation in their sum.
Theorem~\ref{long68:res:prime-pole} tests whether a maximal prime power
survives reduction. Neither supplies the upper bound for a clearing
integer, relative to its remainder, that these arguments need.

\medskip
\noindent No irrationality conclusion is obtained here.  The finite
conclusion is that every rational representation $\Ser=a/q$, $q>0$, must
satisfy both $q\nmid299999!$ and $q\ge2^{39990}>10^{12038}$.
The criteria above require a condition at arbitrarily large indices.
Extending either finite computation alone does not establish it.

\section*{Sources and evidence}

The links below identify formal statements and proofs in the source
records supplied with this revision. They retain their original commits
and paths; they are not a claim that all linked modules were compiled
together at \texttt{\commitshort}.

The attached declaration index is for public snapshot
\nolinkurl{6b78209ab63a8c643281115f8628a3be79ff7ec7} and the separate release
\nolinkurl{52f29ad173b04e3bac941b3663f2b9aebe5de0bb}.
It distinguishes recorded checked declarations from source outside the
recorded build and from release-only declarations. In particular,
\texttt{PaperCompleteExisting} is present but outside the recorded checked
build; the earlier references to its absence are superseded by this source
packet. The same outside-build qualification applies to the candidate
finite-size certificate.

The build receipt refers to a successful Lean build at an earlier commit;
the Lean build step at the supplied public pin was skipped. No Lean build
or axiom audit was run for this prose revision. A source declaration, a
recorded proof-checking result, and a numerical computation are different
kinds of evidence. Descriptions below of checked statements refer to the
supplied records, not to a new compilation.  Theorem~\ref{long68:res:lcm-growth} and
Corollary~\ref{long68:res:radius-constant} are ordinary proofs given in full above,
with corresponding Lean source statements
\href{https://github.com/wcook04/plectis-erdos/blob/0b500c7cf8e8bb7ae343484378df02f277fb8194/lean/ErdosProblems/Erdos68/PaperCompleteLiminf.lean#L42}{the liminf bound for the common denominator} and
\href{https://github.com/wcook04/plectis-erdos/blob/0b500c7cf8e8bb7ae343484378df02f277fb8194/lean/ErdosProblems/Erdos68/PaperCompleteLiminf.lean#L54}{the asymptotic lower bound for the support parameter}.  The finite-block inequality is the separate Lean source
\href{https://github.com/wcook04/plectis-erdos/blob/d788dd4b8c59f2246000f2ed98fffb8a5e8ac72e/lean/ErdosProblems/Erdos68/ChannelIntegralCongruence.lean#L524-L529}{the inequality for a terminal block}.  The two prefix cancellations, the index-$52$
example and the continued-fraction enclosure are finite integer calculations
with the procedures displayed above.  The carry census through $300000$ is a
separate exact-interval computation with the receipt named in
\S\ref{long68:sec:finite}.

\begingroup\small
\begin{longtable}{@{}p{.46\linewidth}p{.48\linewidth}@{}}
\toprule
\papertablehead
Statement & Formal counterpart\\
\midrule
\endfirsthead
\toprule
Statement & Formal counterpart (continued)\\
\midrule
\endhead
\bottomrule
\endfoot
Theorem~\ref{long68:res:prime-pole} &
\lword{Erdos68/PrimePoleCriterion.lean}{223}{factorialGapPrefix_den_full_prime_power_iff}{maximal-power survival},
with the residue formula at
\lword{Erdos68/PrimePoleCriterion.lean}{129}{factorialGapPrefixLCMNumerator_mod_prime}{line 129}
and the denominator form at
\lword{Erdos68/PrimePoleDenominator.lean}{41}{natRatio_den_factorization_eq_iff}{line 41}\\
Proposition~\ref{long68:res:wilson-cofinality} &
\lword{Erdos68/PrimeZeroBranch.lean}{3097}{cofinal_prefixPrivate_factorialGap_hits}{first prime occurrences at arbitrarily large indices};
reflection at
\lword{Erdos68/PrimeZeroBranch.lean}{3139}{prime_dvd_reflected_factorialGap_of_odd}{line 3139}\\
Theorem~\ref{long68:res:carry-equivalence} &
\lword{Erdos68/FactorialGapPlateauCore.lean}{986}{irrational_factorialGapSeries_iff_cofinal_strictFacTopRat_misses}{failures of divisibility for the next integer};
denominator exclusions at
\lword{Erdos68/FactorialGapPlateauCore.lean}{812}{rational_denominator_ge_of_nonunit_carry}{line 812}
and
\lword{Erdos68/FactorialGapPlateauCore.lean}{836}{rational_denominator_not_dvd_pred_factorial_of_nonunit_carry}{line 836}\\
Proposition~\ref{long68:res:companion-orbit} &
\lword{Erdos68/CompanionOrbitRationality.lean}{438}{not_irrational_factorialGapSeries_iff_eventually_companion_floor_neg_two}{factorial residues and rationality}\\
Proposition~\ref{long68:res:lower-escape} &
\lword{Erdos68/ShrinkingTargetNormalForm.lean}{82}{irrational_factorialGapSeries_iff_cofinal_lower_endpoint_escape}{the equivalent fractional-part inequality};
finite window at
\lword{Erdos68/ShrinkingTargetNormalForm.lean}{122}{lower_endpoint_escape_of_predecessorGap_outside_window}{line 122},
irrationality implication at
\lword{Erdos68/ShrinkingTargetNormalForm.lean}{145}{irrational_factorialGapSeries_of_cofinal_predecessorGap_outside_window}{line 145},
classification \eqref{long68:eq:escape-classification} at
\lword{Erdos68/ShrinkingTargetNormalForm.lean}{39}{factorialGap_lower_endpoint_escape_iff_nonunit_or_maximal_digit}{line 39},
radius at
\lword{Erdos68/ShrinkingTargetNormalForm.lean}{158}{lower_endpoint_target_radius_lt_two_div_sq}{line 158}\\
Theorem~\ref{long68:res:shift-family} &
\lword{Erdos68/FactorialShiftFamilyOrbit.lean}{139}{not_irrational_shiftGapSeries_iff_eventually_ceil_residue_two}{the criterion for the shifted series};
escape form at
\lword{Erdos68/FactorialShiftFamilyOrbit.lean}{156}{irrational_shiftGapSeries_iff_cofinal_ceil_residue_misses}{line 156},
member $t=-1$ at
\lword{Erdos68/FactorialShiftFamilyOrbit.lean}{175}{shiftGapSeries_neg_one_eq_factorialGapSeries}{line 175},
and the irrationality of $e$ at
\lword{Erdos68/FactorialShiftFamilyOrbit.lean}{211}{irrational_exp_one_of_family_boundary}{line 211}\\
Theorem~\ref{long68:res:global-residue} &
\lword{Erdos68/EndpointWeightedPrivateSupport.lean}{6164}{irrational_factorialGapSeries_of_cofinal_global_complementaryTail}{the prime-parameter specialization};
the natural-parameter wrapper is linked after the proof and is outside the
recorded checked build. Coprimality \eqref{long68:eq:gap-normalisation} is at
\lword{Erdos68/EndpointWeightedPrivateSupport.lean}{4209}{factorialBlockTailNumerator_coprime_privateModulus}{line 4209}\\
Equation~\eqref{long68:eq:factor-split} &
\lword{Erdos68/EndpointWeightedPrivateSupport.lean}{4766}{factorialBlock_unitFactorPairFloor_scale_iff}{the two simultaneous inequalities},
floor identity at
\lword{Erdos68/EndpointWeightedPrivateSupport.lean}{4701}{factorialBlock_unitFactorPairFloor_eq_min}{line 4701},
irrationality implication at
\lword{Erdos68/PrimeZeroBranch.lean}{1574}{irrational_factorialGapSeries_of_cofinal_largePrefixPrivate_unitScaleSplit}{line 1574}\\
Removal of a fixed denominator factor &
\lword{Erdos68/EndpointWeightedPrivateSupport.lean}{2677}{factorialBlockPrivateQuotient_eq_one_of_gap_lt}{a fixed factor eventually divides the factorial}\\
Weight identities used in Theorem~\ref{long68:res:normalform} and
\eqref{long68:eq:channel-congruence} &
\lword{Erdos68/FactorialChannelCertificate.lean}{25}{factorial_pow_floor_dvd_factorial}{integrality of the weights},
\lword{Erdos68/FactorialChannelCertificate.lean}{43}{channelWeight_mul_denominator}{the denominator-times-weight identity}\\
Theorem~\ref{long68:res:channel-radius} &
\lword{Erdos68/ChannelIntegralCongruence.lean}{1084}{square_subsequence_radius_three_halves_lower}{finite radius bound};
sequence form at
\lword{Erdos68/ChannelIntegralCongruence.lean}{1103}{no_eventual_square_subsequence_three_halves_upper}{line 1103},
little-$o$ form at
\lword{Erdos68/ChannelIntegralCongruence.lean}{905}{not_isLittleO_square_subsequence_radius}{line 905},
limitation of the numerical inequalities proved here
(the attached release also contains
\nolinkurl{sharp_radius_satisfies_square_log_constraint};
this is separate from the historical link above)\\
Changing one weighted sum $V_p$ &
\lword{Erdos68/PrimeUnitTranslator.lean}{1657}{exists_remote_factorialGrid_primeTranslator_reduction}{adjustment at sufficiently large indices},
residual identity at
\lword{Erdos68/PrimeUnitTranslator.lean}{1559}{exists_cramerFactorialGridResidual_eq_det_mul_factorialGapSeries_add_int}{line 1559}\\
Equation~\eqref{long68:eq:doubled-prime} &
\lword{Erdos68/FactorialZeroPlateauSupplement.lean}{257}{sq_dvd_double_strictFacTop_factorialGapPrefix_iff}{doubled-prime criterion}\\
Canonical factorial digits &
\lword{Erdos68/CanonicalFactorialTermination.lean}{78}{exists_rat_eq_iff_eventually_zero_canonicalDigit}{termination equivalence}\\
Amplification modulus &
\lword{Erdos68/PrimeZeroBranch.lean}{5557}{factorialGapRepeatedBirthAmplificationModulus_dvd_succ_den}{divisibility},
nonvanishing at
\lword{Erdos68/PrimeZeroBranch.lean}{5572}{predecessorGapNumeratorNat_mod_repeatedBirthAmplificationModulus_ne_zero}{line 5572}\\
\end{longtable}
\endgroup

The two weight identities in the table are ingredients, not the complete
congruence statement. In the supplied public snapshot,
\nolinkurl{ChannelIntegralCongruence.lean} contains
\nolinkurl{channelNumerator_mod_factorialMoment} (line 76) and
\nolinkurl{exists_channelCorrection} (line 127). The former states
$(d!-1)\mid M(c)-V_d(c)$; negating the integer quotient gives the sign
convention of Theorem~\ref{long68:res:normalform}. These are source
correspondences, not a fresh build or axiom audit.

Attribution.  Wilson's theorem and the Wilson reflection identity are
classical, the latter recorded by Stewart \cite[p.~462, (4)]{stewart2004}.
The factorial-digit termination criterion goes back to Cantor
\cite{cantor1869}; Koepf and Schmersau prove its irrationality direction for
digits that are not eventually maximal
\cite[Example~3.2, p.~121]{koepf-schmersau}, and Galambos treats rationality
criteria for Cantor series \cite[Ch.~II, \S2.1]{galambos1976}.  Here the identity
$C=\Ser-e+2$ gives the eventual digit value $m-2$ and extends to the shifted
family. The comparison with Han\v{c}l and Tijdeman
\cite[Lemma~2.1 and the following
remark, p.~385]{hancl-tijdeman} concerns factorial scaling: their lemma
makes a scaled tail integral, whereas the proof in
\S\ref{long68:sec:carry} identifies the next integer above a generally
nonintegral scaled prefix.  The multiplicity bound is Garaev,
Luca and Shparlinski's \cite{garaev-luca-shparlinski}, and the lcm deduction
from it is not theirs.  The divisibility in Lemma~\ref{long68:res:gap-gcd} is
the case $P=-1$ of the relation used in the proof of Lemma~5 of Luca and
Shparlinski \cite{luca-shparlinski} and at display~(2.5) of Lai \cite{lai}.  The
nonvanishing cutoff for polynomial shifts is Lemma~3 of
\cite{luca-shparlinski}, restated with the bound $n!+P(n)>1$ in
\cite[Lemma~2.1]{lai}.  The survival criterion is a specialisation of
Louwsma and Martino's valuation formula \cite[Lemma~4.1, p.~10]{louwsma-martino}.
The single-denominator growth criterion is Erd\H{o}s's
\cite[Theorem~1, p.~1]{erdos1975}, and the continued-fraction identities are
those of \cite[\S1.12(ii)]{nist-dlmf}. The deductions from Wilson's
theorem, the conditional finite support bound and
Theorem~\ref{long68:res:lcm-growth} are proved above. No further priority
claim is inferred from this comparison.


\appendix
\section{Guide to the formal sources}\label{long68:app:sources}

The following links identify the statements used above and related
lemmas. Each preserves the original file, declaration, and line reference.
The evidence qualifications in the preceding section apply throughout;
in particular, source presence alone does not establish inclusion in a
checked build.

\subsection*{Formalisation coverage and remaining dependencies}
\label{long68:sec:coverage}
\papersectiontarget{coverage}
Every theorem, lemma, proposition and corollary of this record has a Lean
statement of the same assertion, with the same hypotheses, checked by the
Lean kernel using only the axioms \texttt{propext}, \texttt{Classical.choice}
and \texttt{Quot.sound}, in the development at revision
\texttt{181078b6b009}.  No result of this record depends on an input that is
absent from that development.  The concordance after the list below gives
the declarations for each result.

\begin{itemize}[leftmargin=*]
\item \lword{Erdos68/CanonicalFactorialDigits.lean}{47}{facFloor_rat_eq_cleared}{the factorial-scaled floor of a rational number}
\item \lword{Erdos68/CanonicalFactorialDigits.lean}{81}{canonicalDigit_eq_zero_of_rational}{eventual zero digits of a rational number}
\item \lword{Erdos68/CanonicalFactorialDigits.lean}{124}{canonicalDigit_eq_floor_mul_remainder}{a digit as the floor of a scaled fractional part}
\item \lword{Erdos68/CanonicalFactorialDigits.lean}{153}{canonicalDigit_lt_radix}{the digit is smaller than its radix}
\item \lword{Erdos68/CanonicalFactorialDigits.lean}{162}{canonicalRemainder_recurrence}{the recurrence for fractional parts}
\item \lword{Erdos68/CanonicalFactorialDigits.lean}{196}{factorial_expansion_partial}{the finite factorial expansion with its remainder}
\item \lword{Erdos68/CanonicalFactorialDigits.lean}{269}{zero_remainder_forces_zero_digit_tail}{a zero remainder forces all later digits to vanish}
\item \lword{Erdos68/CanonicalFactorialTermination.lean}{38}{canonicalRemainder_eq_zero_of_zero_digit_tail}{eventual zero digits force a zero remainder}
\item \lword{Erdos68/CanonicalFactorialTermination.lean}{54}{exists_rat_eq_of_canonicalRemainder_eq_zero}{a zero factorial remainder implies rationality}
\item \lword{Erdos68/FactorialGapPlateauCore.lean}{82}{factorialGrid_plateau}{the scaled partial sum and limit select the same integer}
\item \lword{Erdos68/FactorialGapPlateauCore.lean}{945}{irrational_factorialGapSeries_of_cofinal_nonunit_carries}{non-unit carries at arbitrarily large indices imply irrationality}
\item \lword{Erdos68/FactorialZeroPlateauCertificates.lean}{136}{sixty_seven_le_rational_denominator}{every rational denominator is at least sixty-seven}
\item \lword{Erdos68/PrimeZeroBranch.lean}{6099}{factorialGap_zero_branch_iff_lower_endpoint_cylinder}{the exact interval condition for the lower endpoint}
\item \lword{Erdos68/PrimeZeroBranch.lean}{3047}{exists_late_prefixPrivate_factorialGap_hit_of_primeProduct_lt}{a prime-product bound gives a late first occurrence}
\item \lword{Erdos68/PrimeZeroBranch.lean}{3231}{prime_dvd_factorialBlockNormalizedCollisionCore_of_odd_reflection}{a reflected prime survives normalization under the block and upper-half hypotheses}
\item \lword{Erdos68/PrimeZeroBranch.lean}{3252}{one_lt_factorialBlockPrimeHitCount_of_odd_reflection}{two distinct in-block indices under the same reflection hypotheses}
\item \lword{Erdos68/PrimeZeroBranch.lean}{1419}{factorialGapLargePrefixPrivatePowerModulus_dvd_tailoredBlockPrivateQuotient}{the selected prime powers divide the remaining quotient}
\item \lword{Erdos68/PrimeZeroBranch.lean}{1524}{factorialGapLargePrefixPrivatePrimes_unit_factor_pair_tailoredBlock}{the factors one and a selected prime in the chosen block}
\item \lword{Erdos68/AdjacentUnitCarryWindow.lean}{148}{consecutive_unit_carries_iff_positive_offset_le_den}{the exact integer interval for two consecutive unit carries}
\item \lword{Erdos68/AdjacentUnitCarryWindow.lean}{215}{twoStep_den_mul_transitionNormalizers}{the two denominator-reduction factors telescope}
\item \lword{Erdos68/AdjacentUnitCarryWindow.lean}{243}{adjacentUnitCarryWindowDen_eq_twoStep_den}{the denominator after two steps}
\item \lword{Erdos68/AdjacentUnitCarryWindow.lean}{337}{adjacentUnitCarryWindowOffset_eq_twoStep_reducedNumerator}{the offset equals the later numerator times the reduction factors}
\item \lword{Erdos68/ChannelBreakpointRigidity.lean}{71}{factorialMoment_eq_factorial_pow_mul_channelNumerator_band}{factorisation when the floor in the weights is constant}
\item \lword{Erdos68/ChannelBreakpointRigidity.lean}{91}{factorialMoment_eq_zero_of_channelNumerator_eq_zero_band}{vanishing of the moment on one such interval}
\item \lword{Erdos68/ChannelBreakpointRigidity.lean}{101}{factorialMoment_eq_factorial_mul_channelNumerator_firstBand}{factorisation on the interval from d to twice d}
\item \lword{Erdos68/ChannelBreakpointRigidity.lean}{130}{exists_index_ge_two_mul_of_factorialMoment_ne_zero_of_channel_eq_zero}{a nonzero moment forces an index at least twice d}
\item \lword{Erdos68/ChannelIntegralCongruence.lean}{868}{square_subsequence_radius_cubic_lower}{a cubic lower bound for the support parameter}
\item \lword{Erdos68/ChannelIntegralCongruence.lean}{885}{no_eventual_square_subsequence_cubic_upper}{exclusion of the stated eventual cubic upper bound}
\item \lword{Erdos68/FactorialChannelCertificate.lean}{65}{channelEvent_eq_zero_of_not_dvd}{no change unless the divisor divides the index}
\item \lword{Erdos68/FactorialChannelCertificate.lean}{156}{lambda34_residual_bounds}{the strict rational bounds for the example remainder}
\item \lword{Erdos68/FactorialAnalyticBoundary.lean}{52}{tendsto_factorialGap_succ_div_sq}{the consecutive-term quadratic ratio tends to zero}
\item \lword{Erdos68/FactorialAnalyticBoundary.lean}{84}{not_summable_abs_factorialGap_succ_div_sq_sub_one}{divergence of the absolute quadratic-ratio errors}
\item \lword{Erdos68/FactorialAnalyticBoundary.lean}{131}{tendsto_factorialGap_rpow_one_div_two_pow}{the iterated-root growth expression tends to one}
\item \lword{Erdos68/EndpointWeightedPrivateSupport.lean}{556}{pairwiseCollisionCore_insert_gcd_lcm}{updating the pairwise gcd lcm after insertion}
\item \lword{Erdos68/EndpointWeightedPrivateSupport.lean}{586}{collisionCore_insert_gcd_lcm}{updating the lcm with a fixed prescribed factor}
\item \lword{Erdos68/EndpointWeightedPrivateSupport.lean}{698}{collisionCore_div_base_eq_pairwiseCollisionCore_div_gcd}{exact removal of the common part with the base}
\item \lword{Erdos68/EndpointWeightedPrivateSupport.lean}{739}{collisionCore_div_base_factorization}{the valuation after removal of the base}
\item \lword{Erdos68/EndpointWeightedPrivateSupport.lean}{800}{collisionCore_div_base_mul_denominatorLcm_dvd_denominatorProd}{the shared part times the lcm divides the product}
\item \lword{Erdos68/EndpointWeightedPrivateSupport.lean}{2707}{factorialBlockFixedPairPrivateQuotients_eq_one}{fixed denominator indices are eventually absorbed by the factorial}
\item \lword{Erdos68/EndpointWeightedPrivateSupport.lean}{3327}{factorialBlockNormalizedCollisionCore_le_gapProd_div_gapLcm}{the shared part is bounded by the product divided by the lcm}
\item \lword{Erdos68/EndpointWeightedPrivateSupport.lean}{3394}{factorialBlockNormalizedCollisionCore_factorization}{the surviving valuation after removal of the factorial base}
\item \lword{Erdos68/EndpointWeightedPrivateSupport.lean}{3443}{factorialBlock_primePower_le_gapPow_of_two_hits}{a shared prime power is bounded by the gap product}
\item \lword{Erdos68/EndpointWeightedPrivateSupport.lean}{3472}{factorialBlock_prime_hitPair_distance_gt_exponent}{two indices sharing a prime power are farther apart than its exponent}
\item \lword{Erdos68/EndpointWeightedPrivateSupport.lean}{3602}{factorialBlock_primePowerHitCount_mul_succ_le}{the resulting bound on the number of such indices}
\item \lword{Erdos68/EndpointWeightedPrivateSupport.lean}{3790}{factorialBlock_normalizedCollision_factorization_add_base_lt_prime}{the total relevant exponent is less than the prime}
\item \lword{Erdos68/EndpointWeightedPrivateSupport.lean}{3848}{factorialBlock_normalizedCollision_factorization_add_base_lt_blockDiameter}{the total relevant exponent is less than the block diameter}
\item \lword{Erdos68/EndpointWeightedPrivateSupport.lean}{3982}{factorialBlock_prime_sq_gt_pred_of_dvd_normalizedCollisionCore}{a surviving prime has square greater than the parameter minus one}
\item \lword{Erdos68/EndpointWeightedPrivateSupport.lean}{4387}{factorialBlock_factorProjection_lcm_eq_privateModulus}{the projection moduli have the required lcm}
\item \lword{Erdos68/EndpointWeightedPrivateSupport.lean}{4854}{factorialBlock_complementaryFactorPairFloor_le_global}{the two-factor bound is no larger than the full residue bound}
\item \lword{Erdos68/EndpointWeightedPrivateSupport.lean}{5213}{factorialBlock_normalizedCollisionCore_factorization_eq_repeatedHitLayerCount}{the valuation counts exponents occurring at two indices}
\item \lword{Erdos68/EndpointWeightedPrivateSupport.lean}{5277}{factorialBlock_normalizedCollisionCore_factorization_eq_truncatedRepeatedHitLayerCount}{a finite exponent range suffices for that count}
\item \lword{Erdos68/EndpointWeightedPrivateSupport.lean}{5389}{factorialBlock_primePower_dvd_normalizedCollisionCore_iff_one_lt_hitCount_of_endpoint_lt_base}{the exact repeated-power criterion above the factorial base}
\end{itemize}

% , - part extended_record , -
% Long-only authored mathematics for the Erdős #68 reasoning record.
% Everything here was in the short note before the September 2026 revision and
% is kept because the complete record owns the subsidiary arguments, the
% superseded deductions, and the failed routes with their exact boundaries.

% BEGIN GENERATED CONCORDANCE (claude_lane/build_concordance.py; do not edit by hand)
\paragraph{Concordance of statements and Lean declarations.}
Each result of this record that has a kernel-checked Lean statement of the same assertion is listed
below with the declarations that jointly state it.  Each name links to its declaration at revision
\texttt{181078b6b009}.  Where the Lean statement is stronger than the printed one and implies it by
an immediate specialisation, the entry says so.

\begingroup\small\raggedright
\par\noindent\hangindent=1.5em Theorem~\ref{long68:res:prime-pole}: \href{https://github.com/wcook04/plectis-erdos/blob/181078b6b009d905809cf2e007ad309386a3d1e0/lean/ErdosProblems/Erdos68/PaperCompletePrimePole.lean\#L117}{\nolinkurl{maximal_prime_power_survival}}.
\par\noindent\hangindent=1.5em Proposition~\ref{long68:res:wilson-cofinality}: \href{https://github.com/wcook04/plectis-erdos/blob/181078b6b009d905809cf2e007ad309386a3d1e0/lean/ErdosProblems/Erdos68/PaperCompleteExisting.lean\#L180}{\nolinkurl{cofinal_first_prime_occurrences}}.
\par\noindent\hangindent=1.5em Lemma~\ref{long68:res:product-lcm}: \href{https://github.com/wcook04/plectis-erdos/blob/181078b6b009d905809cf2e007ad309386a3d1e0/lean/ErdosProblems/Erdos68/PaperCompleteExisting.lean\#L189}{\nolinkurl{product_lcm_pairwise_gcd}}.
\par\noindent\hangindent=1.5em Lemma~\ref{long68:res:gap-gcd}: \href{https://github.com/wcook04/plectis-erdos/blob/181078b6b009d905809cf2e007ad309386a3d1e0/lean/ErdosProblems/Erdos68/ChannelIntegralCongruence.lean\#L288}{\nolinkurl{factorial_gap_gcd_exact}}.
\par\noindent\hangindent=1.5em Lemma~\ref{long68:res:segment} (the Lean statement is stronger): \href{https://github.com/wcook04/plectis-erdos/blob/181078b6b009d905809cf2e007ad309386a3d1e0/lean/ErdosProblems/Erdos68/ChannelIntegralCongruence.lean\#L566}{\nolinkurl{factorialGapSegment_log_sum_le_channelLCM_add_choose}}.
\par\noindent\hangindent=1.5em Theorem~\ref{long68:res:lcm-growth}: \href{https://github.com/wcook04/plectis-erdos/blob/181078b6b009d905809cf2e007ad309386a3d1e0/lean/ErdosProblems/Erdos68/PaperCompleteLiminf.lean\#L42}{\nolinkurl{common_denominator_growth_liminf}}.
\par\noindent\hangindent=1.5em Theorem~\ref{long68:res:carry-equivalence}: \href{https://github.com/wcook04/plectis-erdos/blob/181078b6b009d905809cf2e007ad309386a3d1e0/lean/ErdosProblems/Erdos68/PaperCompleteExisting.lean\#L77}{\nolinkurl{strict_successor_characterisation}}.
\par\noindent\hangindent=1.5em Proposition~\ref{long68:res:companion-orbit}: \href{https://github.com/wcook04/plectis-erdos/blob/181078b6b009d905809cf2e007ad309386a3d1e0/lean/ErdosProblems/Erdos68/PaperCompleteExisting.lean\#L38}{\nolinkurl{companion_orbit}}.
\par\noindent\hangindent=1.5em Proposition~\ref{long68:res:lower-escape}: \href{https://github.com/wcook04/plectis-erdos/blob/181078b6b009d905809cf2e007ad309386a3d1e0/lean/ErdosProblems/Erdos68/PaperCompleteExisting.lean\#L102}{\nolinkurl{lower_interval_criterion}}.
\par\noindent\hangindent=1.5em Theorem~\ref{long68:res:shift-family}: \href{https://github.com/wcook04/plectis-erdos/blob/181078b6b009d905809cf2e007ad309386a3d1e0/lean/ErdosProblems/Erdos68/PaperCompleteExisting.lean\#L123}{\nolinkurl{uniform_family_boundary}}, \href{https://github.com/wcook04/plectis-erdos/blob/181078b6b009d905809cf2e007ad309386a3d1e0/lean/ErdosProblems/Erdos68/PaperCompleteExisting.lean\#L134}{\nolinkurl{uniform_family_members}}.
\par\noindent\hangindent=1.5em Theorem~\ref{long68:res:global-residue}: \href{https://github.com/wcook04/plectis-erdos/blob/181078b6b009d905809cf2e007ad309386a3d1e0/lean/ErdosProblems/Erdos68/PaperCompleteExisting.lean\#L169}{\nolinkurl{global_complementary_criterion_prime}}.
\par\noindent\hangindent=1.5em Theorem~\ref{long68:res:normalform}: \href{https://github.com/wcook04/plectis-erdos/blob/181078b6b009d905809cf2e007ad309386a3d1e0/lean/ErdosProblems/Erdos68/PaperCompleteSupportedBands.lean\#L17}{\nolinkurl{supported_integral_normal_form}}.
\par\noindent\hangindent=1.5em Theorem~\ref{long68:res:bandbreakpoint}: \href{https://github.com/wcook04/plectis-erdos/blob/181078b6b009d905809cf2e007ad309386a3d1e0/lean/ErdosProblems/Erdos68/PaperCompleteSupportedBands.lean\#L24}{\nolinkurl{supported_quotient_band}}, \href{https://github.com/wcook04/plectis-erdos/blob/181078b6b009d905809cf2e007ad309386a3d1e0/lean/ErdosProblems/Erdos68/PaperCompleteSupportedBands.lean\#L42}{\nolinkurl{supported_first_band_cancellation}}, \href{https://github.com/wcook04/plectis-erdos/blob/181078b6b009d905809cf2e007ad309386a3d1e0/lean/ErdosProblems/Erdos68/PaperCompleteSupportedBands.lean\#L51}{\nolinkurl{supported_breakpoint_escape}}.
\par\noindent\hangindent=1.5em Theorem~\ref{long68:res:moment-ideal}: \href{https://github.com/wcook04/plectis-erdos/blob/181078b6b009d905809cf2e007ad309386a3d1e0/lean/ErdosProblems/Erdos68/PaperCompleteMomentIdeal.lean\#L242}{\nolinkurl{exact_moment_ideal_with_primitive_attainment}}, \href{https://github.com/wcook04/plectis-erdos/blob/181078b6b009d905809cf2e007ad309386a3d1e0/lean/ErdosProblems/Erdos68/PaperCompleteMomentIdeal.lean\#L321}{\nolinkurl{minimumMoment_independent_prime}}.
\par\noindent\hangindent=1.5em Theorem~\ref{long68:res:residual-transparency}: \href{https://github.com/wcook04/plectis-erdos/blob/181078b6b009d905809cf2e007ad309386a3d1e0/lean/ErdosProblems/Erdos68/PaperCompleteResidualIdentity.lean\#L116}{\nolinkurl{residual_transparency}}, \href{https://github.com/wcook04/plectis-erdos/blob/181078b6b009d905809cf2e007ad309386a3d1e0/lean/ErdosProblems/Erdos68/PaperCompleteResidualIdentity.lean\#L166}{\nolinkurl{summable_fullResidual}}, \href{https://github.com/wcook04/plectis-erdos/blob/181078b6b009d905809cf2e007ad309386a3d1e0/lean/ErdosProblems/Erdos68/PaperCompleteResidualIdentity.lean\#L185}{\nolinkurl{zero_moment_residual_integral}}, \href{https://github.com/wcook04/plectis-erdos/blob/181078b6b009d905809cf2e007ad309386a3d1e0/lean/ErdosProblems/Erdos68/PaperCompleteResidualIdentity.lean\#L195}{\nolinkurl{equal_moment_residual_integer_difference}}.
\par\noindent\hangindent=1.5em Theorem~\ref{long68:res:channel-radius}: \href{https://github.com/wcook04/plectis-erdos/blob/181078b6b009d905809cf2e007ad309386a3d1e0/lean/ErdosProblems/Erdos68/PaperCompleteExisting.lean\#L234}{\nolinkurl{square_subsequence_radius}}, \href{https://github.com/wcook04/plectis-erdos/blob/181078b6b009d905809cf2e007ad309386a3d1e0/lean/ErdosProblems/Erdos68/PaperCompleteExisting.lean\#L258}{\nolinkurl{radius_no_eventual_ratio_upper}}, \href{https://github.com/wcook04/plectis-erdos/blob/181078b6b009d905809cf2e007ad309386a3d1e0/lean/ErdosProblems/Erdos68/PaperCompleteExisting.lean\#L277}{\nolinkurl{radius_not_littleO}}.
\par\noindent\hangindent=1.5em Corollary~\ref{long68:res:radius-constant}: \href{https://github.com/wcook04/plectis-erdos/blob/181078b6b009d905809cf2e007ad309386a3d1e0/lean/ErdosProblems/Erdos68/PaperCompleteLiminf.lean\#L54}{\nolinkurl{asymptotic_radius_constant_liminf}}.
\par\noindent\hangindent=1.5em Theorem~\ref{long68:res:translator}: \href{https://github.com/wcook04/plectis-erdos/blob/181078b6b009d905809cf2e007ad309386a3d1e0/lean/ErdosProblems/Erdos68/PaperCompleteExisting.lean\#L299}{\nolinkurl{prime_channel_corrector}}.
\par\endgroup
% END GENERATED CONCORDANCE

\section{Further deductions and limitations}\label{long68:sec:extended-record}

This appendix supplies details used by the earlier sections: factorial
digits and their remainders, exact gcd and lcm identities, and additional
finite examples. It also records weaker deductions and explains why several
proposed irrationality arguments do not establish their needed hypotheses.

\subsection{Factorial digits and their remainders}\label{long68:sec:ext-digits}

For a real number $x$, write $\theta_0=\{x\}$ and, for $m\ge1$,
\[
 a_m=\lfloor m\,\theta_{m-1}\rfloor,\qquad \theta_m=m\,\theta_{m-1}-a_m ,
\]
so that $\theta_m=\{m!\,x\}$ for $m\ge1$.  The kernel checks the floor
formula, the digit bounds $0\le a_m<m$, the recurrence
$\theta_{m+1}=(m+1)\theta_m-a_{m+1}$, the finite telescoping expansion
\[
 x=\lfloor x\rfloor+\sum_{m=2}^{N}\frac{a_m}{m!}+\frac{\theta_N}{N!},
\]
and the rule that a zero remainder at one index forces every later digit to
vanish.  The rational direction is also checked: if $a\in\mathbb Z$, $q\in\mathbb N$, and $0<q\le n$, then
$\lfloor n!a/q\rfloor=(n!/q)a$ and the canonical digit at
radix $n+1$ vanishes, so every rational input has an eventually zero
expansion.  The converse holds for every real input, and gives Cantor's
termination criterion \cite{cantor1869}
\[
 x\in\Q\quad\Longleftrightarrow\quad a_m(x)=0\ \hbox{for all large }m .
\]
If all digits after index $N\ge1$ vanish, the recurrence gives
$\theta_{N+k}\ge(k+1)\theta_N$; since every remainder is below one,
$\theta_N=0$ and $x=\lfloor N!x\rfloor/N!$.

For $x=S$, a zero canonical digit means $m\theta_{m-1}<1$. The
stronger condition $m\theta_{m-1}<E_m$, where $E_m=m!(S-H_m)$, is the
lower-interval event used in \S\ref{long68:sec:carry}. The reported interval
data contain zero digits at $m=5$ and $m=23$, but no occurrence of that
stronger condition through $m=100000$.

There is also a useful identity for an abstract rational sequence $F_m$
and integer sequence $k_m$. Suppose
$F_m=mF_{m-1}+1+\varepsilon_m-k_m$, and define
\[
 \delta_m=\lceil F_m\rceil-F_m,\qquad
 q_m=m\lceil F_{m-1}\rceil+1-k_m-\lceil F_m\rceil.
\]
Then $0\le\delta_m<1$, and the exact identities are
\[
 q_m=\lfloor m\delta_{m-1}-\varepsilon_m\rfloor,\qquad
 \delta_m=m\delta_{m-1}-\varepsilon_m-q_m.
\]
The attached Lean source proves these identities under the displayed
recurrence. They do apply to the actual partial sums: take $F_m=m!H_m$,
$k_m=0$ and $\varepsilon_m=1/(m!-1)$ for $m\ge3$. The recurrence then
follows from $H_m=H_{m-1}+1/(m!-1)$, starting with $F_2=2$.
The ceiling code $q_m$ need not equal the carry $b_m$, because $Z_m$ is
the least integer strictly above $F_m$, not always its ceiling. Precisely,
$Z_m=\lceil F_m\rceil+\mathbf1_{\{F_m\in\Z\}}$, so
\[
 b_m=q_m+m\mathbf1_{\{F_{m-1}\in\Z\}}
          -\mathbf1_{\{F_m\in\Z\}}.
\]
For example, $F_2=2$ and $F_3=36/5$ give $q_3=-1$ but $b_3=2$.
The codes agree whenever both scaled prefixes are nonintegral. Once
$F_2=2$ and $\varepsilon_m=1/(m!-1)$ are fixed, the recurrence with $k_m=0$
determines every $F_m$ uniquely by induction. Thus matching the recurrence
and its initial value is stronger than satisfying selected congruences
or carry bounds. This specialization is an ordinary deduction from the
partial-sum recurrence, not a new Lean check. It does not show that the rational gaps admit a
finite-state description or establish the unbounded non-unit carries
needed for irrationality.

\subsection{How the classical criteria apply here}\label{long68:sec:ext-literature}

Let $s_n<a$ be partial sums tending to $a$. Koepf and Schmersau show that
$\lfloor ns_n\rfloor=\lfloor na\rfloor$ for all sufficiently large $n$
forces $a$ to be irrational \cite[Thm.~1.1, p.~117]{koepf-schmersau}.
The strict inequality matters: if $a=A/B$ were rational, then every
sufficiently large multiple $n$ of $B$ would give
$\lfloor ns_n\rfloor<na=\lfloor na\rfloor$. Their rational-term version
uses a positive integer multiplier $p_n$ and obtains the floor equality
from $np_ns_n\in\mathbb N_0$ and the strict tail bound
$a-s_n<1/(np_n)$
\cite[(2.1) and (2.3), p.~118; Thms.~2.2--2.3, pp.~119--120]{koepf-schmersau}.
Indeed, integrality makes $ns_n$ a multiple of $1/p_n$, so its gap to the
least strictly larger integer is at least $1/p_n$; the increase
$n(a-s_n)<1/p_n$ cannot cross that integer.
For the choice they record after (2.1), the least common multiple of the
reduced summand denominators, here $p_n=\lcm\{k!-1:2\le k\le n\}$, the last
two denominators already obstruct the tail bound. As proved in
\S\ref{long68:sec:lcm}, they are coprime for $n\ge3$, so
\[
 p_n\ge(n!-1)\bigl((n-1)!-1\bigr).
\]
For $n\ge4$ this gives $np_n>(n+1)!-1$. Hence the first omitted summand
$1/((n+1)!-1)$ already exceeds $1/(np_n)$, and this $p_n$ cannot satisfy
their tail hypothesis. This rules out the displayed choice of $p_n$, not
the criterion itself. In fact, the least positive integer $p_n$ for which
$np_nH_n$ is integral is exactly
\[
 p_n=\den(nH_n)
     =\frac{\den(H_n)}{\gcd(n,\den(H_n))}.
\]
Indeed, writing $H_n$ in lowest terms shows that its denominator must
divide $np_n$. With this least choice the required inequality is
\[
 0<S-H_n<\frac{\gcd(n,\den(H_n))}{n\den(H_n)}.
\]
Thus reduction of the partial sum, and the common factor of its reduced
denominator with $n$, determine the best scale available in this
criterion. The eventual inequality above remains unproved here; failure
of a larger clearing multiplier does not establish its failure.

Duverney's Theorem~3.1 includes the quadratic growth assumption
$cu_n^2\le u_{n+1}\le c'u_n^2$ for positive constants $c,c'$, while
$u_{n+1}/u_n^2\to0$ for $u_n=n!-1$
\cite[(1.3), p.~275; Thm.~3.1, pp.~285--286]{duverney}.
His Corollary~3.2, which allows signs $a_n\in\{-1,1\}$, assumes convergence
of the signed series $\sum_n(u_{n+1}/u_n^2-1)$ in (3.6)
\cite[Cor.~3.2, p.~287]{duverney}. Its terms tend to $-1$ here, so the
series does not converge. The supplied Lean records check the ratio
limit and also the divergence of the absolute-value series; the signed
condition fails already by the term test. To see the ratio limit directly, divide
numerator and denominator by $(n!)^2$: the numerator tends to zero and
the denominator $(1-1/n!)^2$ tends to one.

For a strictly increasing sequence $(n_k)$ of positive integers,
Erd\H{o}s proved irrationality of $\sum_k1/n_k$ under the following
conditions, with a fixed $\varepsilon>0$ \cite[Theorem~1, p.~1]{erdos1975}:
\[
 \limsup_{k\to\infty}n_k^{1/2^k}=\infty,
 \qquad n_k>k^{1+\varepsilon}\quad\text{eventually}.
\]
Barreto, Kang, Kim, Kova\v{c}
and Zhang treat products of consecutive denominators and weighted extensions
\cite[Thms.~2--3 and Rem.~4, pp.~2--5]{barreto-et-al}.  For $a_n=n!-1$ and every fixed $\psi>1$,
\[
 0\le\frac{\log(n!-1)}{\psi^n}\le\frac{n^2}{\psi^n}\longrightarrow0,
 \qquad (n!-1)^{1/\psi^n}\longrightarrow1.
\]
Thus neither Erd\H{o}s's limsup hypothesis, which uses $\psi=2$, nor the
growth hypotheses of their Theorems~2 and~3 hold. The supplied Lean limit
is the special case $\psi=2$; the estimate above also covers the other
bases greater than one used in those theorems.  Their proof of Theorem~3 uses the classical
criterion, which they trace to Fourier's proof that $e$ is irrational, that a
rational sum of nonnegative rationals with infinitely many positive terms
admits no prefix-clearing integers $D_N$ with $\liminf_ND_Nr_N=0$, where $r_N$
is the tail after $N$ terms \cite[Lem.~8, p.~6]{barreto-et-al}; their
Proposition~12 produces such integers under the hypotheses of that theorem
\cite[pp.~9--12]{barreto-et-al}.
For the present series, even the full summand lcm $\Lc_N$ makes
$\Lc_N(S-H_N)$ tend to infinity, as shown in \S\ref{long68:sec:lcm}.
A smaller integer cannot clear every summand. It can nevertheless clear
the partial sum after addition: the least positive such integer is
$\den(H_N)$, and every other one is a multiple of it. Thus this criterion
would require information about cancellation in the reduced partial sum,
not merely an improvement from the product to the lcm. More precisely, if
$S=a/q$, then
\[
 q\den(H_N)(S-H_N)\in\mathbb Z_{>0},
 \qquad \den(H_N)(S-H_N)\ge\frac1q.
\]
Thus $\liminf_N\den(H_N)(S-H_N)=0$ is a sufficient target for this
prefix-clearing criterion; the preceding common-denominator estimates
do not establish it.

Dividing the recurrence $Z_m=mZ_{m-1}+1-b_m$ by $m!$ and
telescoping gives the exact finite identity
\[
 \frac{Z_M}{M!}=\frac{Z_2}{2!}+\sum_{m=3}^{M}\frac{1-b_m}{m!},
\]
so the carry defects $1-b_m$ are integer coefficients of a factorial series.
Han\v{c}l and Tijdeman classify the rational sums with polynomial coefficients
\cite[Thm.~3.1 and Cor.~3.1, pp.~390--391]{hancl-tijdeman}; their
denominator is the cumulative linear product $\prod_{n\le N}(an+b)$, and the
individual number $N!-1$ does not occur.  In the factorial case $a=1$, $b=0$, the criterion of Oppenheim that they
reproduce \cite[Lem.~2.2, p.~385]{hancl-tijdeman} says the following: for
integer coefficients $c_n$ with $|c_n|<n$ eventually and
$\liminf_n|c_n|/n=0$, the sum $\sum_n c_n/n!$ is rational exactly when
$c_n=0$ eventually. This formulation concerns integer coefficients and
factorial denominators, not reciprocal terms with denominator $n!-1$.
However, their introduction records a version of Oppenheim's result
without the liminf assumption: under $|c_m|<m-1$ eventually, rationality
is equivalent to eventual vanishing
\cite[Introduction, p.~383]{hancl-tijdeman}. Changing finitely many terms
adds a rational number, so eventual hypotheses suffice. Here
$2-m\le c_m=1-b_m\le2$ gives $|c_m|<m-1$ for $m\ge4$; the classical
result therefore applies directly.

For completeness, the following elementary proof specialises that result
to these one-sided bounds. For $N\ge2$, telescoping
$\sum_{m>N}(m-1)/m!=1/N!$ gives
\[
 -1<N!\sum_{m>N}\frac{c_m}{m!}
 \le2N!\sum_{m>N}\frac1{m!}<\frac2N.
\]
The lower bound follows from $c_m\ge2-m$ and
$N!\sum_{m>N}(m-2)/m!=1-N!\sum_{m>N}1/m!<1$.
The upper bound follows by comparison with
$\sum_{j\ge1}(N+1)^{-j}=1/N$.
If $\sum_{m\ge3}c_m/m!$ is rational, the scaled tails are integers for all
sufficiently large $N$. They must then be zero, and subtracting consecutive
tails gives $c_m=0$ eventually. The converse is immediate. Since
$Z_M/M!\to S$, the displayed finite identity recovers
Theorem~\ref{long68:res:carry-equivalence}. This is a self-contained
specialisation of the classical rationality criterion, not a strengthening
of it. It proves an equivalence, not the nonvanishing needed for irrationality.

\subsection{A superseded deduction}\label{long68:sec:ext-superseded}

A multiplicity theorem gives a weaker lcm bound, which is useful for
comparison with the elementary proof in \S\ref{long68:sec:lcm}.
For an odd prime $r$, let $m_r$ count the
indices $2\le n\le N$ with $r\mid n!-1$; such an index satisfies $n<r$.  The
factorial-congruence multiplicity estimate of Garaev, Luca and Shparlinski
\cite[arXiv v1, Thm.~12, p.~16]{garaev-luca-shparlinski}, applied to the
residue $a=1$ on $1\le n\le\min(N,r-1)$, gives $m_r\ll N^{2/3}$; the prime $2$
divides none of
these factors. The sum of the exponents at a prime is at most the number
of occurrences times the largest exponent, which is its exponent in
$\Lc_N$. Thus
\[
 \sum_{n=2}^{N}\log(n!-1)=\sum_r\sum_{n=2}^{N}v_r(n!-1)\log r
 \le\bigl(\max_r m_r\bigr)\log\Lc_N
 \ll N^{2/3}\log \Lc_N ,
\]
and Stirling's formula makes the left side $\asymp N^2\log N$, whence
$\log \Lc_N\gg N^{4/3}\log N$.  Theorem~\ref{long68:res:lcm-growth} supersedes this:
its exponent is $3/2$, its constant is explicit, and it does not use the
external multiplicity theorem.  The deduction is retained because it is the only place a
multiplicity bound enters the record.

Dividing all coefficients by their gcd does not evade the lcm restriction:
the resulting integer vector still has $V_2=\cdots=V_D=0$, so its moment
is still divisible by $L_D$. For the nonzero cofactor vector in
\S\ref{long68:sec:open}, normalization is justified by the elementary
argument given there. It supplies no proof that the remainder is
nonintegral. No formalised primitive cofactor construction is claimed here.

\subsection{Rational grid points and first crossings}
\label{long68:sec:ext-plateau}

One can instead compare a partial sum $H$ with the rational numbers
having a fixed positive denominator $q$. The next such number above $H$
is $(\lfloor qH\rfloor+1)/q$. It is at most $\Ser$ exactly when
$\lfloor qH\rfloor+1\le q\Ser$.  Suppose $H<G\le S$, the number $n!G$ is an integer, and
$n!(S-H)<1$. Then
\[
 n!H<n!G\le n!S<n!H+1.
\]
Thus $n!G$ is both the least integer strictly above $n!H$ and the floor of
$n!S$. In particular, for $n\ge2$, suppose a single rational level $G$
satisfies $H_{n+1}<G\le S$ and $n!G\in\Z$. Since $H_n<H_{n+1}$,
the tail bounds at both indices give
\[
 \lfloor(n+1)!S\rfloor=(n+1)!G
   =(n+1)\lfloor n!S\rfloor.
\]
The next canonical digit is therefore zero. The same level must persist
at both indices; choosing an unrelated rational level at each index
would not imply this equality.

Set $H_1=0$, the empty partial sum. Now let $\tau\ge2$ be a first
crossing of a rational level $G$, so that
$H_{\tau-1}<G\le H_\tau$, and write $G-H_{\tau-1}=a/v$ with positive
integers $a,v$. The newly added summand gives
\[
 0<\frac av\le\frac1{\tau!-1},
 \qquad v\ge a(\tau!-1)\ge\tau!-1.
\]
The scaled overshoot is
\[
 0\le\tau!(H_\tau-G)
 =1+\frac1{\tau!-1}-\tau!\frac av<2.
\]
Consequently $-\lfloor\tau!(H_\tau-G)\rfloor$ is $0$ or $-1$.
Here the minus sign is outside the floor; taking the floor of the negative
overshoot would give a different integer at nonintegral overshoots.
The displayed integer is $-1$ precisely when the overshoot is at least
$1$, or equivalently when
\[
 \frac av\le\frac1{\tau!(\tau!-1)}.
\]
In that case $v\ge\tau!(\tau!-1)$. This extra small-gap condition is not
asserted for every crossing. Nor is $-\lfloor\tau!(H_\tau-G)\rfloor$
automatically the carry $b_\tau$ defined from $Z_\tau$: here $G$ is an
arbitrary rational level. Neither denominator bound requires $a/v$ to be in lowest terms.

There is also a direct obstruction at prime indices. With $\Delta_m$ as
defined in \S\ref{long68:sec:carry}, the formal source proves for $m\ge3$
the criterion
\[
 m\mid Z_m
 \quad\Longleftrightarrow\quad
 1+\frac1{m!-1}<m\Delta_m\le2+\frac1{m!-1},
\]
and if $\Ser=a/q$ with $q>0$, then $p\mid Z_p$ for every prime $p>q$, so one
exact missed prime $p$ gives $q\ge p$.  Exact kernel reduction gives
$11\nmid Z_{11}$, hence $q\ge11$; exact rational normalisation gives
$60\nmid Z_{60}$, $64\nmid Z_{64}$ and $67\nmid Z_{67}$, and the prime index
$67$ gives the checked bound
\[
 \Ser=\frac aq,\ q>0\quad\Longrightarrow\quad q\ge67 .
\]
The exact-interval census of the short note replaces $67$ by $300000$.

The finite geometric-series identity gives another exact decomposition.  For a real $x\ne0,1$ and an integer $K\ge0$,
\[
 \frac1{x-1}=\sum_{j=1}^{K}\frac1{x^j}+\frac1{x^K(x-1)} .
\]
When $x=k!$ and a chosen factorial scale is divisible by $(k!)^K$, the scaled
finite sum is integral and only the last term retains the factor $k!-1$ in
its denominator.  The identity isolates one residual fraction before exact
bounding, and it supplies no cofinal family of nonzero residuals.

The earlier record rejects a proposed divisibility strengthening of the
first-crossing bound and reports examples at $m=52$ and $m=591$. Its
notation for the rational quantity in that claim was not defined, so those
reports are not used as a verified counterexample here. The established
conclusion remains the lower bound on the denominator's size, not a
specified factor dividing it.

There is a second, more arithmetic mechanism at doubled prime indices, stated
in \eqref{long68:eq:doubled-prime}.  The formal theorem is not
restricted to individually computed indices; the divisibility criterion
holds for the actual partial sums at every odd prime.

\subsection{Exact identities for shared prime powers}\label{long68:sec:ext-collision}

The following identities make it possible to update the shared part of
a common denominator and to calculate the prime powers left after removal
of a factorial factor. Let $I\subseteq\{2,3,\ldots\}$ be finite.

Write $D(I)=\lcm_{i<j,\ i,j\in I}\gcd(d_i,d_j)$, with $D(I)=1$
for an empty or singleton family. For each prime $r$, its valuation is the
second-largest valuation among the $d_i$, counting missing values as zero;
the denominator lcm has the largest valuation. This proves
$D(I)\lcm_{i\in I}d_i\mid\prod_{i\in I}d_i$ prime by prime.

The integer $D(I)$ can be updated exactly when one index is added. For the positive
factorial-gap denominators, adjoining a new index $a\ge2$, $a\notin I$, gives
\[
 D(I\cup\{a\})=\lcm\!\Bigl(D(I),\ \gcd\bigl(d_a,\lcm_{j\in I}d_j\bigr)\Bigr),
\]
since finite-family gcd and lcm distributivity collapses the lcm of all
pairwise gcds against $d_a$ to a single gcd. For a fixed positive integer
$F$, take the lcm with $F$ on both sides:
\[
 \lcm\bigl(F,D(I\cup\{a\})\bigr)
 =\lcm\Bigl(F,D(I),\gcd\bigl(d_a,\lcm_{j\in I}d_j\bigr)\Bigr).
\]
Thus one need only retain the denominator lcm and $\lcm(F,D(I))$ at each
step. The prescribed factor $F$ is not an extra summand denominator.

There is also an exact bound by the product of the denominators divided
by their least common multiple. For a positive integer $F$, define
$\widetilde C(I)=\lcm(F,D(I))/F=D(I)/\gcd(F,D(I))$. Since
$\widetilde C(I)$ divides $D(I)$, the preceding divisibility gives
\[
 \widetilde C(I)\lcm_{j\in I}d_j\mid\prod_{j\in I}d_j,\qquad
 \widetilde C(I)\le\frac{\prod_{j\in I}d_j}{\lcm_{j\in I}d_j}.
\]  For the factorial block this specialises to
\[
 \widetilde C_p\le\frac{\prod_{n\in I_p}(n!-1)}{\lcm_{n\in I_p}(n!-1)},
\]
an upper bound for the shared factor in terms of the denominator product
and lcm. The bound alone does not establish
\eqref{long68:eq:weighted-target}.

Dividing by $F_p$ subtracts its prime exponents; it does not remove every
prime that occurs in $F_p$. With $F_p=(p-1)!$,
$D_p=D(I_p)$ and $C_p=\lcm(F_p,D_p)$, as in the main record,
\[
 \widetilde C_p=\frac{C_p}{F_p}=\frac{D_p}{\gcd(F_p,D_p)},\qquad
 v_r(\widetilde C_p)=\max\{0,v_r(D_p)-v_r(F_p)\} .
\]
For $e>0$, $r^e\mid\widetilde C_p$ exactly when $D_p$ is divisible by
$r^{e+v_r(F_p)}$, which in the block forces two distinct denominators to be divisible
by that higher power, by the second-largest-valuation formula. If $i<j$
are two such hits and $f=v_r(F_p)$, then
$r^{e+f}\mid j!/i!-1$. Since $r\mid j!-1$ forces $j<r$,
$0<j!/i!-1<r^{j-i}$, hence $e+f<j-i$.
For a prime $r\mid\widetilde C_p$, this bounds the surviving valuation by
$v_r(\widetilde C_p)+v_r(F_p)<r$. Since $v_r(\widetilde C_p)\ge1$ and
$v_r(F_p)\ge\lfloor(p-1)/r\rfloor$, we obtain
$\lfloor(p-1)/r\rfloor\le r-2$, hence $p-1<r(r-1)<r^2$.
Consequently $\widetilde C_p$ is coprime to $k!$ whenever
$k(k-1)\le p-1$: a prime $r\le k$ dividing both would give
$p-1<r(r-1)\le k(k-1)\le p-1$. This excludes the small primes but does not bound the
product of the remaining prime powers.

If a prime $r$ divides a denominator $n!-1$ with $n\ge p$,
then $r\nmid F_p$ and $r^e\mid\widetilde C_p\iff r^e\mid C_p$ for every
$e>0$.  This has an exact incidence-count form,
\[
 r^e\mid\widetilde C_p
 \quad\Longleftrightarrow\quad
 1<\#\{i\in I_p:r^e\mid i!-1\},
\]
so a bound of at most one such index implies $v_r(\widetilde C_p)<e$.
Counting the exponents gives
\[
 v_r(\widetilde C_p)=\#\bigl\{e\in[1,r-1]:1<\#\{i\in I_p:r^e\mid i!-1\}\bigr\},
\]
and for a prime $r>2p-1$ the range can be restricted to
$e\in[1,2p-4]$.  The same equivalence holds under the exact condition
$r\nmid F_p$ in place of $r>2p-1$, hence in particular for every prime
$r\ge p$. At $e=2$, an upper bound of one on the number of indices gives
$v_r(\widetilde C_p)\le1$.

A surviving prime power also forces two indices to be far apart.  If $r$ is prime, $e>0$ and
$r^e\mid\widetilde C_p$, there are $i<j$ in $I_p$ with
$r^{e+v_r(F_p)}\mid i!-1$, $r^{e+v_r(F_p)}\mid j!-1$ and
$r^{e+v_r(F_p)}\le j^{\,j-i}$; consequently $(2p-1)^d<r^{e+v_r(F_p)}$ forces
$d<j-i$.  The required spacing follows directly: any two $r^e$-hits
$i<j$ satisfy $e<j-i$, because $r\mid j!-1$ already forces $j<r$ and the
inequality for $j!/i!-1$ then forces the strict separation. More generally,
for any two indices $2\le i<j$ and any prime power $r^e$ with $e>0$,
\[
 r^e\mid i!-1,\quad r^e\mid j!-1
 \quad\Longrightarrow\quad r^e\le j^{\,j-i},
\]
so successive solutions of $i!\equiv1\pmod{r^e}$ are at least $e+1$
apart, and
\[
 (e+1)\,\#\{i\in I_p:r^e\mid i!-1\}\ \le\ 2p+e-2 .
\]
The resulting bound on the exponent is
\[
 r\mid\widetilde C_p
 \quad\Longrightarrow\quad
 v_r(\widetilde C_p)+v_r(F_p)<2p-3
\]
for every prime $r$ and $p\ge2$. Thus the exponent already present in
$F_p$ uses part of the same bound $2p-3$. To apply
\eqref{long68:eq:global-scale}, one still needs sufficiently strong bounds
for the product of the shared prime powers and for the gap $\rho_p/R_p$.

The available squarefreeness evidence is finite. An exhaustive modular scan through
$r\le2{,}000{,}000$ and $n\le240$ found four individual square hits and no
prime with two such hits.  All $498{,}501$ pairs $2\le a<b\le1000$ have
squarefree $\gcd(a!-1,b!-1)$. The supplied aggregate scan through $p=499$ also reports a ratio below
$0.374$, but does not specify its logarithmic normalisation; that ratio is
not used as a quantitative premise here. These are reported finite
computations, not an asymptotic bound or a proof of squarefreeness at all
indices.

A finite version of the argument for first prime occurrences compares the product of a
chosen set of primes, each at least $5$, with $\prod_{2\le k\le B}(k!-1)$; if
the prime product is larger, at least one chosen prime has no hit through
$B$, while Wilson still bounds its least hit by $q-2$.  Wilson reflection
limits what a linear-size divisor can be assumed to be: if $n$ is odd,
$n<q$, and $q\mid n!-1$, then $q\mid(q-n-1)!-1$. Suppose also that
$p\le n$, both indices lie in $I_p$, and the reflected index is earlier,
equivalently $q<2n+1$. The prime $q$ then divides two denominators in the
block. Since $q>n\ge p$, it does not divide $F_p=(p-1)!$, so
$q\mid\widetilde C_p$: removing the factorial factor does not remove this
shared prime.

Stewart states that for every $\varepsilon>0$ there are infinitely many odd
$n$ whose least prime factor $q$ of $n!-1$ satisfies
\[
 n<q<\left(\frac{\sqrt{145}-1}{8}+\varepsilon\right)n ,
\]
estimate~(9) of Theorem~1 being stated for $n!+1$ \cite[p.~463]{stewart2004}
and transferred to $n!-1$ in the text \cite[p.~464]{stewart2004}.  This controls $q$ relative to the specified index $n$, not to its first
occurrence. It does not by itself force a repeated divisor. For the minus
sign, the bound is already met at all sufficiently large Wilson indices
$n=q-2$. For every prime $q\ge5$, Wilson gives $q\mid(q-2)!-1$, and $q$ is its least prime
factor. Indeed, every prime divisor exceeds $q-2$, and $q-1$ is even;
moreover $q/(q-2)\to1$. The reflected index is then $1$, outside the denominators of $S$.

For example, modulo $11$ the only solution of $k!\equiv1$ with
$2\le k\le9$ is $k=9$, although $11/9<(\sqrt{145}-1)/8$.
A small ratio $q/n$ therefore supplies neither a second admissible hit nor
membership of a reflected index in the chosen block. The earlier
reflection argument needs both block-membership hypotheses and distinct
indices. Neither the required estimate at first occurrences nor a bound
for the total shared part follows from the transferred least-prime-factor
estimate alone.

For a selected prime $q\mid R_p$, the factors $1$ and $q$ are coprime,
and the associated moduli $R_p$ and $R_p/q$ have lcm $R_p$. Thus this
specialisation requires no second selected prime.

The elementary fact behind the projection argument is as follows. Let
$Z,T,B$ be nonnegative integers, let $R>0$, and suppose $Z\equiv T\pmod R$.
Every positive divisor $Q$ of $R$ with $Z\le B<Q$ satisfies
$T\bmod Q=Z$, since $Z$ is already the least nonnegative representative.
Consequently, unequal residues modulo two divisors exceeding $B$ rule out
such a $Z$. More generally, if $T\bmod Q_1\ne T\bmod Q_2$, then
$\min(Q_1,Q_2)\le T$: otherwise both residues would equal $T$.

\subsection{Finite vectors and exact numerical examples}\label{long68:sec:ext-certificates}

The minimum-moment vector $12K_4+253U_6-11U_8$ from
\S\ref{long68:sec:channels} also gives a finite nonintegral remainder.
By Theorem~\ref{long68:res:residual-transparency}, its remainder is
$242+1380(S-H_4)$. Exact rational arithmetic gives
\[
 255+\frac45
 <242+1380\sum_{d=5}^{8}\frac1{d!-1}
 <255+\frac56,
 \qquad \frac{2\cdot1380}{9!-1}<\frac1{100}.
\]
The tail bound therefore places the remainder strictly between $255$ and
$256$. Its fractional part, rather than the size of the remainder itself,
is what excludes denominators dividing $1380$. This small instance explains
the signed-part comparison in the short note; the headline finite
exclusions already imply its denominator restriction.

The finite-support vector $\lambda=2e_3-e_4$ has, by kernel check,
$\chn2\lambda=0$, factorial moment $-12$, $\chn3\lambda=-2$,
$\chn4\lambda=11$, and $\chn d\lambda=-12$ for every $d\ge5$.  Under the
exact tail enclosure
$1/119<\sum_{d\ge5}1/(d!-1)<1/50$ its residual lies
strictly between $-93/575$ and $-309/13685$, so it is nonzero and has
absolute value less than $1$.

In the enlarged coefficient space of \S\ref{long68:sec:channels}, exact
integer computation verifies the vectors $K_D$ for every $2\le D\le12$:
weighted sums $2$ through $D$ vanish, the factorial moment is $\Lc_D$, and the
coefficients have gcd one.  They lie outside the space of vectors supported
on $n\ge2$, since $\Lc_D$ is odd and every such vector has even moment.  At
$D=9$,
\[
 \Lc_9=31540008254514077395,\qquad a_9=[e_1]K_9=-3902884074990939115 .
\]
Since $U_{11}=T_{11}$ has $u_{11}=0$, the vector $K_9-9553024718754\,U_{11}$
keeps the coordinate $a_9$ and has coefficients with gcd one.  By
Theorem~\ref{long68:res:residual-transparency} its residual is
$\Lc_9(\Ser-H_9)-9553024718754$, and exact rational arithmetic with the tail
bound $\sum_{n\ge36}1/(n!-1)<2/(36!-1)$ places it strictly between
$1353/100000$ and $1354/100000$.  A different example is supported on
$n\ge2$: the vector $c=(-40,55,-10,1)$ on the support $(3,4,5,6)$ annihilates
weighted sums $2$ and $3$, has moment $600$, and satisfies
\[
 0.09925341997208298<\mathcal R(c)<0.09925341997208300 ,
\]
which excludes denominators dividing $600$.

A separate interval computation reports the stronger geometric statement that no
lower-interval event $m\theta_{m-1}<E_m$ occurs at any $3\le m\le100000$; its executable and source
digest are not available, so that classification remains external finite
evidence.  The carry census through $300000$ used in the short note is the GMP
computation reported in the supplied manuscript, not a run repeated for
this prose revision. Section~\ref{long68:sec:finite} records its reported
driver, backend and payload digests and the separate earlier calculation
through $4000$.  The full range is external
computational evidence outside the Lean development. The original prose
packet contains the manuscript's report, not the complete GMP executable
and payload, so the present revision does not independently authenticate
those digests.

\subsection{Limits of the recurrences and clearing factors used here}\label{long68:sec:ext-nogo}

\begin{itemize}
\item Without the defining floor relation, the carry recurrence and its
  range allow $Z_m/m!$ to be the constant $3/2$. Taking $Z_m=3m!/2$ for
  $m\ge2$ and $b_m=1$ for $m\ge3$ gives
  $Z_m=mZ_{m-1}+1-b_m$ and $-1\le b_m\le m-1$, with the actual initial
  value $Z_2=3$. But it gives $Z_3=9$, whereas the actual prefix
  $H_3=6/5$ gives $Z_3=8$. This example only disproves sufficiency of
  the stated recurrence, bounds and initial integer. It fails the floor
  definition. No claim is made that it satisfies the additional prime-index
  identities. The rational recurrence for $F_m=m!H_m$ in
  \S\ref{long68:sec:ext-digits} determines the actual partial sums uniquely.
\item The arguments considered with Wilson quotients, harmonic sums,
  $p$-adic gamma identities, and factorial residues still need a real
  gap estimate and the required modular divisibility at the same indices.
  No such unbounded family is obtained here. This records the missing
  step in these arguments, not an impossibility theorem for the use of
  those identities.
\item For the genus-zero product $E(z)=\prod_{n\ge2}(1-z/n!)$, local uniform
  convergence and logarithmic differentiation give $-E'(1)/E(1)=\Ser$.
  Termwise clearing by $\prod_{n=2}^N(n!-1)$ cannot give a positive
  remainder tending to zero: this product is at least $L_N$, and
  $L_N(\Ser-H_N)\to\infty$ by \S\ref{long68:sec:lcm}. The conclusion concerns
  this clearing factor alone; Hermite--Pad\'e systems with other
  denominators are not ruled out.
\item Changing a coefficient vector without changing its moment changes
  the residual by an integer only. Also, cancelling the first $D-1$ weighted
  sums forces $L_D$ to divide the moment; factorial divisibility does not
  remove that constraint.
\item A fixed pair of denominator indices cannot make the projection
  argument work at arbitrarily large parameters, because their factors
  eventually divide the factorial being removed.
\end{itemize}

\subsection{An elementary criterion for any real number}
\label{long68:sec:ext-generic}

The fractional-part condition in \eqref{long68:eq:lower-escape} follows
from the following consequence of Cantor's termination criterion.  For a real
$x$ and any positive sequence $c_m$ with $c_m\le1/m$ eventually,
\[
 x\notin\Q
 \quad\Longleftrightarrow\quad
 \{(m-1)!\,x\}\ge c_m\ \hbox{for arbitrarily large }m .
\]
For rational $x$ the fractional parts vanish eventually. Conversely, if
$\{(m-1)!x\}<c_m\le1/m$ for all large $m$, then
$m\{(m-1)!x\}<1$, so every sufficiently late canonical digit is zero.
Cantor's criterion then makes $x$ rational. Equality $c_m=1/m$ is allowed:
it is failure of the displayed weak inequality that makes
$m\{(m-1)!x\}$ strictly less than one.
For $x=S$, the choice $c_m=E_m/m$ gives exactly
\eqref{long68:eq:lower-escape}; indeed $0<E_m<1$ ensures the required
threshold bound. The finite test \eqref{long68:eq:finite-escape} is a
separate sufficient condition at one index. Its cofinal converse uses
the carry theorem, as shown in \S\ref{long68:sec:carry}, not substitution
into this general criterion. Neither argument proves that $S$ meets
the test at arbitrarily large indices.

The threshold $1/m$ cannot be replaced by $\lambda/m$ for a fixed
$\lambda>1$. For $x=e$ and $m\ge3$, the factorial tail gives
\[
 m\{(m-1)!e\}
 =1+m!\sum_{n>m}\frac1{n!}<1+\frac1m<\lambda
 \qquad\text{eventually}.
\]
Thus the irrational number $e$ would never meet the enlarged threshold
at sufficiently large indices. Positivity is also essential: with
$c_m=0$, rational numbers would meet the threshold eventually. These
examples explain the two assumptions without imposing any
equidistribution hypothesis.
% , - part family_catalogue , -
\section{Relations among the criteria}\label{sec:erdos-68-complete-family-map}
The two finite exclusions in \S\ref{long68:sec:finite} use only the carry
criterion and a rational enclosure. The coefficient constructions do not
enter either calculation. They instead produce integer linear forms
$\mathcal R(c)=M(c)S+k$, $k\in\Z$. At fixed moment, changing coefficients
changes only $k$; primitive cofactors on the stated progression reproduce
the explicit progression vector rather than a new family.

For a rational value $S=a/q$, such a form is integral whenever $q\mid M(c)$.
To extend this method beyond finite exclusions, one therefore needs
nonintegral remainders with moments covering every positive denominator
in this divisibility sense. The common-denominator lower bound and the
prime-power survival test do not supply the needed upper bound for a
remainder relative to its integer gap. Likewise, the equivalent carry,
digit and interval criteria still require events at arbitrarily large
indices. A finite calculation can rule out denominators without proving
any of those unbounded assertions.

% , - part back , -
\section*{Acknowledgements}
The author thanks Wouter van Doorn for advice on exposition: explaining notation when it first appears, avoiding private terminology, and saying how restrictive a conditional hypothesis is. His advice concerned the writing of another note; he has not reviewed the mathematics of this paper.

\begin{thebibliography}{99}

\bibitem{erdos1988}
Paul Erd{\H o}s. \href{https://doi.org/10.1017/CBO9780511897184.009}{On the irrationality of certain series: problems and results}.
In Alan Baker (ed.), \emph{New Advances in Transcendence Theory}, Cambridge University Press (1988), pp.~102--109.

\bibitem{bloom}
Thomas F. Bloom. \href{https://www.erdosproblems.com/68}{Erd{\H o}s Problem {\#68}}.
Online resource (2026). Historical access: 28 July 2026; present-page status not reverified.

\bibitem{erdos1975}
Paul Erd{\H o}s. \href{https://www.renyi.hu/~p_erdos/1976-44.pdf}{Some problems and results on the irrationality of the sum of infinite series}.
\emph{Journal of Mathematical Sciences} \textbf{10} (1975), 1--7.

\bibitem{louwsma-martino}
Joel Louwsma and Joseph Martino. \href{https://arxiv.org/abs/2309.07280v1}{Rational numbers with odd greedy expansion of fixed length}.
Preprint (2023). arXiv:2309.07280.

\bibitem{cantor1869}
Georg Cantor. {\"U}ber die einfachen Zahlensysteme.
\emph{Zeitschrift f{\"u}r Mathematik und Physik} \textbf{14} (1869), 121--128.

\bibitem{galambos1976}
J{\'a}nos Galambos. \href{https://doi.org/10.1007/BFb0081642}{Representations of Real Numbers by Infinite Series}.
Lecture Notes in Mathematics 502, Springer (1976).

\bibitem{garaev-luca-shparlinski}
Moubariz Z. Garaev, Florian Luca and Igor E. Shparlinski. \href{https://doi.org/10.1090/S0002-9947-04-03612-8}{Character sums and congruences with {$n!$}}.
\emph{Transactions of the American Mathematical Society} \textbf{356} (12) (2004), 5089--5102. arXiv:math/0403422.

\bibitem{stewart2004}
Cameron L. Stewart. \href{https://doi.org/10.5486/PMD.2004.3190}{On the greatest and least prime factors of {$n!+1$}, {II}}.
\emph{Publicationes Mathematicae Debrecen} \textbf{65} (3--4) (2004), 461--480.

\bibitem{koepf-schmersau}
Wolfram Koepf and Dieter Schmersau. \href{https://doi.org/10.1524/anly.2011.1094}{Irrationality of certain infinite series {II}}.
\emph{Analysis} \textbf{31} (2011), 117--124.

\bibitem{duverney}
Daniel Duverney. \href{https://www.ms.u-tokyo.ac.jp/journal/pdf/jms080206.pdf}{Irrationality of fast converging series of rational numbers}.
\emph{Journal of Mathematical Sciences, the University of Tokyo} \textbf{8} (2001), 275--316.

\bibitem{hancl-tijdeman}
Jaroslav Han{\v c}l and Robert Tijdeman. \href{https://doi.org/10.4064/aa118-4-5}{On the irrationality of factorial series}.
\emph{Acta Arithmetica} \textbf{118} (4) (2005), 383--401.

\bibitem{barreto-et-al}
Kevin Barreto, Jiwon Kang, Sang-hyun Kim, Vjekoslav Kova{\v c} and Shengtong Zhang. \href{https://arxiv.org/abs/2601.21442v3}{Irrationality of rapidly converging series: a problem of {Erd{\H o}s} and {Graham}}.
Preprint (2026). arXiv:2601.21442.

\bibitem{luca-shparlinski}
Florian Luca and Igor E. Shparlinski. \href{https://doi.org/10.1112/S0024609305004923}{Prime divisors of shifted factorials}.
\emph{Bulletin of the London Mathematical Society} \textbf{37} (6) (2005), 809--817.

\bibitem{lai}
Li Lai. \href{https://doi.org/10.1017/S0004972725100543}{On the largest prime divisor of {$n!+1$}}.
\emph{Bulletin of the Australian Mathematical Society} \textbf{113} (3) (2026), 390--403. arXiv:2103.14894.

\bibitem{lai-lupu-sprang}
Li Lai, Cezar Lupu and Johannes Sprang. \href{https://doi.org/10.1007/s40687-025-00559-x}{On the irrationality of certain {$p$}-adic zeta values}.
\emph{Research in the Mathematical Sciences} \textbf{12} (4) (2025), article 77. arXiv:2505.23088.

\bibitem{lai-2adic}
Li Lai. \href{https://doi.org/10.1142/S1793042125500113}{On the irrationality of certain {$2$}-adic zeta values}.
\emph{International Journal of Number Theory} \textbf{21} (1) (2025), 207--235. arXiv:2304.00816.

\bibitem{nist-dlmf}
{NIST Digital Library of Mathematical Functions}. \href{https://dlmf.nist.gov/1.12}{Continued fractions: convergents}.
Online resource (2026).

\bibitem{erdos-stewart1976}
Paul Erd{\H o}s and Cameron L. Stewart. \href{https://doi.org/10.1112/jlms/s2-13.3.513}{On the greatest and least prime factors of {$n!+1$}}.
\emph{Journal of the London Mathematical Society} \textbf{13} (3) (1976), 513--519.

\bibitem{luca-shparlinski-exp}
Florian Luca and Igor E. Shparlinski. \href{https://doi.org/10.5802/jtnb.524}{On the largest prime factor of {$n!+2^n-1$}}.
\emph{Journal de Th{\'e}orie des Nombres de Bordeaux} \textbf{17} (3) (2005), 859--870.

\bibitem{garaev-luca-shparlinski-harmonic}
Moubariz Z. Garaev, Florian Luca and Igor E. Shparlinski. \href{https://doi.org/10.1007/s00209-006-0939-5}{Distribution of harmonic sums and {Bernoulli} polynomials modulo a prime}.
\emph{Mathematische Zeitschrift} \textbf{253} (2006), 855--865.

\bibitem{sondow2006}
Jonathan Sondow. \href{https://arxiv.org/abs/0704.1282v2}{A geometric proof that {$e$} is irrational and a new measure of its irrationality}.
\emph{American Mathematical Monthly} \textbf{113} (7) (2006), 637--641. arXiv:0704.1282. Addendum: \emph{Amer. Math. Monthly} \textbf{114} (2007), 659; reading copy v2.

\bibitem{hu2026}
Xiyu Hu. \href{https://arxiv.org/abs/2608.01781v1}{Factorial residues modulo a prime: beyond the square-root bound}.
Preprint (2026). arXiv:2608.01781. Version 1, 3 August 2026.

\bibitem{garaev-pardo2026}
Moubariz Z. Garaev and Julio C. Pardo. \href{https://doi.org/10.1090/proc/17753}{Additive congruences with factorials modulo a prime}.
\emph{Proceedings of the American Mathematical Society} \textbf{154} (10) (2026), 4179--4190. arXiv:2508.12127. Published online 1 July 2026; assigned October 2026 issue.

\bibitem{hu-september2026}
Xiyu Hu. \href{https://arxiv.org/abs/2609.05652v1}{Lower bounds for some value sets over finite fields: incidence geometry and {Bourgain}'s group expansion theorem}.
Preprint (2026). arXiv:2609.05652. Version 1, 4 September 2026.

\bibitem{banks-et-al2005}
William D. Banks, Florian Luca, Igor E. Shparlinski and Henning Stichtenoth. \href{https://journals.tubitak.gov.tr/math/vol29/iss2/6/}{On the Value Set of $n!$ Modulo a Prime}.
\emph{Turkish Journal of Mathematics} \textbf{29} (2) (2005), 169--174.

\bibitem{klurman-munsch2017}
Oleksiy Klurman and Marc Munsch. \href{https://doi.org/10.5802/jtnb.974}{Distribution of factorials modulo $p$}.
\emph{Journal de Théorie des Nombres de Bordeaux} \textbf{29} (1) (2017), 169--177.

\bibitem{grebennikov-et-al2024}
Alexandr Grebennikov, Arsenii Sagdeev, Aliaksei Semchankau and Aliaksei Vasilevskii. \href{https://doi.org/10.4171/RMI/1422}{On the sequence $n!$ mod $p$}.
\emph{Revista Matemática Iberoamericana} \textbf{40} (2) (2024), 637--648.

\bibitem{stevens-dezeeuw2017}
Sophie Stevens and Frank de Zeeuw. \href{https://doi.org/10.1112/blms.12077}{An improved point-line incidence bound over arbitrary fields}.
\emph{Bulletin of the London Mathematical Society} \textbf{49} (5) (2017), 842--858. arXiv:1609.06284.

\bibitem{macleod-barrodale1970}
R. A. Macleod and I. Barrodale. \href{https://doi.org/10.4153/CMB-1970-052-8}{On Equal Products of Consecutive Integers}.
\emph{Canadian Mathematical Bulletin} \textbf{13} (2) (1970), 255--259.

\bibitem{nguyen-tho2026}
{Nguyen Xuan Tho}. \href{https://doi.org/10.4171/EM/556}{On equal products of consecutive integers}.
\emph{Elemente der Mathematik} \textbf{81} (3) (2026), 119--124. First published online 10 July 2025.

\bibitem{erdos-selfridge1975}
Paul Erd{\H o}s and John L. Selfridge. \href{https://doi.org/10.1215/ijm/1256050816}{The product of consecutive integers is never a power}.
\emph{Illinois Journal of Mathematics} \textbf{19} (2) (1975), 292--301.

\bibitem{nist-vandermonde}
{NIST Digital Library of Mathematical Functions}.
\href{https://dlmf.nist.gov/1.3.E13}{Vandermonde determinant},
\S1.3(ii), formula~(1.3.13). Online resource, consulted 18 September 2026.

\end{thebibliography}

% ---- part extended_record ----
% ---- part family_catalogue ----
% ---- part back ----
\end{document}
